% Recovered from the previously completed PDF; no new mathematical content.
% Prose is editable; formula boxes retain original symbol positions.
% Build from repository root: xelatex -no-shell-escape -halt-on-error -output-directory=... print/fudan/NAME.tex
\documentclass[11pt,a4paper]{article}
\usepackage[margin=23mm,headheight=14pt,headsep=8.5mm,footskip=10mm]{geometry}
\usepackage{amsmath,amssymb}
\DeclareMathAlphabet{\mathsf}{OT1}{cmss}{m}{n}
\usepackage[no-math]{fontspec}
\usepackage{xeCJK}
\setmainfont{Latin Modern Roman}
\setCJKmainfont{FandolSong-Regular.otf}[Path=assets/fonts/fandol/,BoldFont=FandolSong-Bold.otf,ItalicFont=FandolSong-Regular.otf,Script=CustomDefault]
\xeCJKsetup{PunctStyle=quanjiao,CheckSingle=true}
\usepackage{setspace}\setstretch{1.12}
\setlength{\parindent}{2em}\setlength{\parskip}{3pt}
\emergencystretch=6em
\hfuzz=9pt
\usepackage{titlesec}
\titleformat{\section}{\centering\normalfont\Large\bfseries}{}{0pt}{}
\titlespacing*{\section}{0pt}{1.8\baselineskip}{1.0\baselineskip}
\usepackage{fancyhdr}\pagestyle{fancy}\fancyhf{}
\fancyhead[L]{\small 数学手记·高等代数}\fancyhead[R]{\small 19A14–19A14、20S01–20S06}
\fancyfoot[C]{\thepage}\renewcommand{\headrulewidth}{0.4pt}
\usepackage[hidelinks,unicode]{hyperref}
\hypersetup{pdftitle={19A14–19A14、20S01–20S06：题目、解答与小结},pdfauthor={Vincent Arqain}}
\newcommand{\RB}[4]{\raisebox{0pt}[#2bp][#3bp]{\hbox to #1bp{#4\hss}}}
\newcommand{\RG}[5]{\hbox to 0pt{\kern#1bp\raisebox{#2bp}[0pt][0pt]{\fontsize{#3bp}{#3bp}\csname font#4\endcsname\char#5}\hss}}
\newcommand{\RT}[4]{\hbox to 0pt{\kern#1bp\raisebox{#2bp}[0pt][0pt]{\fontsize{#3bp}{#3bp}\selectfont#4}\hss}}
\newcommand{\RS}[4]{\hbox to 0pt{\kern#1bp\raisebox{#2bp}[0pt][0pt]{\fontsize{#3bp}{#3bp}\selectfont\ensuremath{#4}}\hss}}
\newcommand{\RR}[4]{\hbox to 0pt{\kern#1bp\raisebox{#2bp}[0pt][0pt]{\rule{#3bp}{#4bp}}\hss}}
\newcommand{\fontmi}{\usefont{OML}{cmm}{m}{it}}
\newcommand{\fontsy}{\usefont{OMS}{cmsy}{m}{n}}
\newcommand{\fontex}{\usefont{OMX}{cmex}{m}{n}}
\newcommand{\fontrm}{\usefont{OT1}{cmr}{m}{n}}
\newcommand{\fontbf}{\usefont{OT1}{cmr}{bx}{n}}
\newcommand{\fontam}{\usefont{U}{msa}{m}{n}}
\newcommand{\fontbb}{\usefont{U}{msb}{m}{n}}
\widowpenalty=10000\clubpenalty=10000\raggedbottom
\begin{document}
\begin{center}{\LARGE\bfseries 19A14–19A14、20S01–20S06\par}\vspace{8pt}复旦大学高等代数习题集·题目与解答\end{center}
% Original page 1.
\section*{2019A14}
设\allowbreak{}\hskip.1em \RB{8.366}{8.312}{0.500}{%
\RG{0.000}{0.001}{10.9091}{mi}{65}%
}\hskip.1em\allowbreak{} 为数域\allowbreak{}\hskip.1em \RB{10.198}{7.952}{0.500}{%
\RG{0.000}{0.001}{10.9091}{mi}{75}%
}\hskip.1em\allowbreak{} 上的\allowbreak{}\hskip.1em \RB{6.729}{5.323}{0.619}{%
\RG{0.000}{0.001}{10.9091}{mi}{110}%
}\hskip.1em\allowbreak{} 阶方阵，满足tr\allowbreak{}\hskip.1em \RB{36.707}{8.683}{3.226}{%
\RG{0.000}{0.001}{10.9091}{rm}{40}\RG{4.243}{0.001}{10.9091}{mi}{65}\RG{12.425}{0.001}{10.9091}{rm}{41}\RG{19.680}{0.001}{10.9091}{rm}{61}\RG{31.189}{0.001}{10.9091}{rm}{48}%
}\hskip.1em\allowbreak{} 。证明：\allowbreak{}\hskip.1em \RB{8.366}{8.312}{0.500}{%
\RG{0.000}{0.001}{10.9091}{mi}{65}%
}\hskip.1em\allowbreak{} 相似于一个主对角元全为\allowbreak{}\hskip.1em \RB{5.518}{7.767}{0.739}{%
\RG{0.000}{0.001}{10.9091}{rm}{48}%
}\hskip.1em\allowbreak{} 的矩阵。

证明：对\allowbreak{}\hskip.1em \RB{6.729}{5.322}{0.620}{%
\RG{0.000}{0.000}{10.9091}{mi}{110}%
}\hskip.1em\allowbreak{} 归纳。\allowbreak{}\hskip.1em \RB{26.147}{7.766}{0.620}{%
\RG{0.000}{0.000}{10.9091}{mi}{110}\RG{9.578}{0.000}{10.9091}{rm}{61}\RG{21.076}{0.000}{10.9091}{rm}{49}%
}\hskip.1em\allowbreak{} 时，\allowbreak{}\hskip.1em \RB{28.229}{8.311}{0.740}{%
\RG{0.000}{0.000}{10.9091}{mi}{65}\RG{11.212}{0.000}{10.9091}{rm}{61}\RG{22.710}{0.000}{10.9091}{rm}{48}%
}\hskip.1em\allowbreak{} ，结论成立。设结论对\allowbreak{}\hskip.1em \RB{24.953}{7.766}{0.620}{%
\RG{0.000}{0.000}{10.9091}{mi}{110}\RG{8.973}{0.000}{10.9091}{sy}{0}\RG{19.882}{0.000}{10.9091}{rm}{49}%
}\hskip.1em\allowbreak{} 阶方阵成立，以下考虑\allowbreak{}\hskip.1em \RB{26.492}{7.766}{1.995}{%
\RG{0.000}{0.000}{10.9091}{mi}{110}\RG{9.579}{0.000}{10.9091}{sy}{21}\RG{21.094}{0.000}{10.9091}{rm}{50}%
}\hskip.1em\allowbreak{} 。

若\allowbreak{}\hskip.1em \RB{39.264}{8.307}{2.220}{%
\RG{0.000}{-0.004}{10.9091}{mi}{65}\RG{11.212}{-0.004}{10.9091}{rm}{61}\RG{22.727}{-0.004}{10.9091}{mi}{21}\RG{29.087}{-0.004}{10.9091}{mi}{73}\RG{33.886}{-1.640}{7.9701}{mi}{110}%
}\hskip.1em\allowbreak{} ，则\allowbreak{}\hskip.1em \RB{24.307}{8.067}{0.646}{%
\RG{0.000}{-0.004}{10.9091}{mi}{110}\RG{6.546}{-0.004}{10.9091}{mi}{21}\RG{15.941}{-0.004}{10.9091}{rm}{61}%
}\hskip.1em\allowbreak{} tr\allowbreak{}\hskip.1em \RB{36.706}{8.678}{3.231}{%
\RG{0.000}{-0.004}{10.9091}{rm}{40}\RG{4.242}{-0.004}{10.9091}{mi}{65}\RG{12.424}{-0.004}{10.9091}{rm}{41}\RG{19.679}{-0.004}{10.9091}{rm}{61}\RG{31.188}{-0.004}{10.9091}{rm}{48}%
}\hskip.1em\allowbreak{} 。由于\allowbreak{}\hskip.1em \RB{10.198}{7.947}{0.504}{%
\RG{0.000}{-0.004}{10.9091}{mi}{75}%
}\hskip.1em\allowbreak{} 是数域，故\allowbreak{}\hskip.1em \RB{26.410}{8.067}{0.744}{%
\RG{0.000}{-0.004}{10.9091}{mi}{21}\RG{9.393}{-0.004}{10.9091}{rm}{61}\RG{20.891}{-0.004}{10.9091}{rm}{48}%
}\hskip.1em\allowbreak{} ，此时无需变换。若\allowbreak{}\hskip.1em \RB{8.366}{8.313}{0.500}{%
\RG{0.000}{0.002}{10.9091}{mi}{65}%
}\hskip.1em\allowbreak{} 不是纯量矩阵，则存在\allowbreak{}\hskip.1em \RB{34.436}{7.976}{0.934}{%
\RG{0.000}{0.002}{10.9091}{mi}{118}\RG{8.710}{0.002}{10.9091}{sy}{50}\RG{19.013}{0.002}{10.9091}{mi}{75}\RG{29.058}{3.961}{7.9701}{mi}{110}%
}\hskip.1em\allowbreak{} ，使\allowbreak{}\hskip.1em \RB{24.293}{8.313}{2.603}{%
\RG{0.000}{0.002}{10.9091}{mi}{118}\RG{5.662}{0.002}{10.9091}{mi}{59}\RG{10.506}{0.002}{10.9091}{mi}{65}\RG{18.687}{0.002}{10.9091}{mi}{118}%
}\hskip.1em\allowbreak{} 线性无关。事实上，若不存在这样的\allowbreak{}\hskip.1em \RB{5.606}{5.324}{0.618}{%
\RG{0.000}{0.002}{10.9091}{mi}{118}%
}\hskip.1em\allowbreak{} ，则每个非零向量都是\allowbreak{}\hskip.1em \RB{8.366}{8.308}{0.503}{%
\RG{0.000}{-0.003}{10.9091}{mi}{65}%
}\hskip.1em\allowbreak{} 的特征向量。取任一组基\allowbreak{}\hskip.1em \RB{44.511}{5.319}{2.608}{%
\RG{0.000}{-0.003}{10.9091}{mi}{101}\RG{5.079}{-1.639}{7.9701}{rm}{49}\RG{9.811}{-0.003}{10.9091}{mi}{59}\RG{14.644}{-0.003}{10.9091}{mi}{58}\RG{19.488}{-0.003}{10.9091}{mi}{58}\RG{24.331}{-0.003}{10.9091}{mi}{58}\RG{29.164}{-0.003}{10.9091}{mi}{59}\RG{34.008}{-0.003}{10.9091}{mi}{101}\RG{39.133}{-1.639}{7.9701}{mi}{110}%
}\hskip.1em\allowbreak{} ，设

\[
\RB{51.989}{8.314}{2.213}{%
\RG{0.000}{0.003}{10.9091}{mi}{65}\RG{8.182}{0.003}{10.9091}{mi}{101}\RG{19.673}{0.003}{10.9091}{rm}{61}\RG{31.188}{0.003}{10.9091}{mi}{21}\RG{40.933}{0.003}{10.9091}{mi}{101}\RG{49.394}{0.003}{10.9091}{mi}{58}\RG{13.261}{-1.633}{7.9701}{mi}{105}\RG{37.552}{-1.633}{7.9701}{mi}{105}\RG{46.012}{-1.633}{7.9701}{mi}{105}%
}
\]
对任意\allowbreak{}\hskip.1em \RB{23.135}{8.310}{2.857}{%
\RG{0.000}{-0.001}{10.9091}{mi}{105}\RG{6.789}{-0.001}{10.9091}{sy}{54}\RG{6.789}{-0.001}{10.9091}{rm}{61}\RG{18.304}{-0.001}{10.9091}{mi}{106}%
}\hskip.1em\allowbreak{} ，向量\allowbreak{}\hskip.1em \RB{30.665}{6.859}{3.763}{%
\RG{0.000}{-0.001}{10.9091}{mi}{101}\RG{5.079}{-1.637}{7.9701}{mi}{105}\RG{10.885}{-0.001}{10.9091}{rm}{43}\RG{21.794}{-0.001}{10.9091}{mi}{101}\RG{26.873}{-1.637}{7.9701}{mi}{106}%
}\hskip.1em\allowbreak{} 也是特征向量，因而

\[
\RB{112.017}{8.677}{3.767}{%
\RG{0.000}{-0.005}{10.9091}{mi}{21}\RG{9.745}{-0.005}{10.9091}{mi}{101}\RG{20.630}{-0.005}{10.9091}{rm}{43}\RG{31.539}{-0.005}{10.9091}{mi}{21}\RG{42.285}{-0.005}{10.9091}{mi}{101}\RG{54.777}{-0.005}{10.9091}{rm}{61}\RG{66.293}{-0.005}{10.9091}{mi}{22}\RG{72.866}{-0.005}{10.9091}{rm}{40}\RG{77.109}{-0.005}{10.9091}{mi}{101}\RG{87.994}{-0.005}{10.9091}{rm}{43}\RG{98.903}{-0.005}{10.9091}{mi}{101}\RG{108.364}{-0.005}{10.9091}{rm}{41}\RG{6.364}{-1.641}{7.9701}{mi}{105}\RG{14.824}{-1.641}{7.9701}{mi}{105}\RG{37.903}{-1.641}{7.9701}{mi}{106}\RG{47.364}{-1.641}{7.9701}{mi}{106}\RG{82.188}{-1.641}{7.9701}{mi}{105}\RG{103.982}{-1.641}{7.9701}{mi}{106}%
}
\]
对某个\allowbreak{}\hskip.1em \RB{30.104}{7.952}{2.855}{%
\RG{0.000}{0.001}{10.9091}{mi}{22}\RG{9.603}{0.001}{10.9091}{sy}{50}\RG{19.906}{0.001}{10.9091}{mi}{75}%
}\hskip.1em\allowbreak{} 成立。比较系数，得\allowbreak{}\hskip.1em \RB{34.446}{8.072}{3.761}{%
\RG{0.000}{0.001}{10.9091}{mi}{21}\RG{6.364}{-1.635}{7.9701}{mi}{105}\RG{12.775}{0.001}{10.9091}{rm}{61}\RG{24.290}{0.001}{10.9091}{mi}{21}\RG{30.654}{-1.635}{7.9701}{mi}{106}%
}\hskip.1em\allowbreak{} ，从而\allowbreak{}\hskip.1em \RB{8.366}{8.312}{0.500}{%
\RG{0.000}{0.001}{10.9091}{mi}{65}%
}\hskip.1em\allowbreak{} 是纯量矩阵，矛盾。

取\allowbreak{}\hskip.1em \RB{29.963}{5.320}{2.139}{%
\RG{0.000}{-0.002}{10.9091}{mi}{101}\RG{5.080}{-1.639}{7.9701}{rm}{49}\RG{12.842}{-0.002}{10.9091}{rm}{61}\RG{24.357}{-0.002}{10.9091}{mi}{118}%
}\hskip.1em\allowbreak{} 、\allowbreak{}\hskip.1em \RB{38.145}{8.309}{2.139}{%
\RG{0.000}{-0.002}{10.9091}{mi}{101}\RG{5.079}{-1.639}{7.9701}{rm}{50}\RG{12.842}{-0.002}{10.9091}{rm}{61}\RG{24.357}{-0.002}{10.9091}{mi}{65}\RG{32.539}{-0.002}{10.9091}{mi}{118}%
}\hskip.1em\allowbreak{} ，并扩张成\allowbreak{}\hskip.1em \RB{15.424}{7.971}{0.502}{%
\RG{0.000}{-0.002}{10.9091}{mi}{75}\RG{10.046}{3.956}{7.9701}{mi}{110}%
}\hskip.1em\allowbreak{} 的一组基。记相应的过渡矩阵为\allowbreak{}\hskip.1em \RB{8.726}{7.949}{0.502}{%
\RG{0.000}{-0.002}{10.9091}{mi}{80}%
}\hskip.1em\allowbreak{} 。由于\allowbreak{}\hskip.1em \RB{41.920}{8.309}{2.139}{%
\RG{0.000}{-0.002}{10.9091}{mi}{65}\RG{8.182}{-0.002}{10.9091}{mi}{101}\RG{13.261}{-1.639}{7.9701}{rm}{49}\RG{21.024}{-0.002}{10.9091}{rm}{61}\RG{32.539}{-0.002}{10.9091}{mi}{101}\RG{37.618}{-1.639}{7.9701}{rm}{50}%
}\hskip.1em\allowbreak{} ，有

\[
\RB{192.694}{19.588}{14.129}{%
\RG{51.084}{18.651}{10.9092}{ex}{32}\RG{85.252}{18.651}{10.9092}{ex}{33}\RG{60.366}{7.276}{10.9091}{rm}{48}\RG{77.331}{7.276}{10.9091}{mi}{11}\RG{8.519}{4.500}{7.9701}{sy}{0}\RG{15.105}{4.500}{7.9701}{rm}{49}\RG{0.000}{-0.004}{10.9091}{mi}{80}\RG{19.837}{-0.004}{10.9091}{mi}{65}\RG{28.019}{-0.004}{10.9091}{mi}{80}\RG{39.568}{-0.004}{10.9091}{rm}{61}\RG{95.706}{-0.004}{10.9091}{mi}{59}\RG{122.357}{-0.004}{10.9091}{mi}{66}\RG{134.225}{-0.004}{10.9091}{sy}{50}\RG{144.528}{-0.004}{10.9091}{mi}{77}\RG{171.568}{-0.004}{10.9091}{rm}{40}\RG{175.811}{-0.004}{10.9091}{mi}{75}\RG{185.857}{-0.004}{10.9091}{rm}{41}\RG{190.100}{-0.004}{10.9091}{mi}{58}\RG{155.111}{-1.640}{7.9701}{mi}{110}\RG{160.250}{-1.640}{7.9701}{sy}{0}\RG{166.836}{-1.640}{7.9701}{rm}{49}\RG{59.720}{-7.899}{10.9091}{mi}{12}\RG{76.422}{-7.899}{10.9091}{mi}{66}%
}
\]
迹在相似变换下不变，故

\[
\RB{156.244}{10.300}{3.231}{%
\RG{61.654}{4.500}{7.9701}{sy}{0}\RG{68.241}{4.500}{7.9701}{rm}{49}\RT{0.000}{-0.004}{10.9091}{t}\RT{4.244}{-0.004}{10.9091}{r}\RG{8.520}{-0.004}{10.9091}{rm}{40}\RG{12.762}{-0.004}{10.9091}{mi}{66}\RG{21.584}{-0.004}{10.9091}{rm}{41}\RG{28.839}{-0.004}{10.9091}{rm}{61}\RT{40.372}{-0.004}{10.9091}{t}\RT{44.616}{-0.004}{10.9091}{r}\RG{48.892}{-0.004}{10.9091}{rm}{40}\RG{53.134}{-0.004}{10.9091}{mi}{80}\RG{72.973}{-0.004}{10.9091}{mi}{65}\RG{81.154}{-0.004}{10.9091}{mi}{80}\RG{89.674}{-0.004}{10.9091}{rm}{41}\RG{96.928}{-0.004}{10.9091}{rm}{61}\RT{108.462}{-0.004}{10.9091}{t}\RT{112.706}{-0.004}{10.9091}{r}\RG{116.982}{-0.004}{10.9091}{rm}{40}\RG{121.224}{-0.004}{10.9091}{mi}{65}\RG{129.406}{-0.004}{10.9091}{rm}{41}\RG{136.661}{-0.004}{10.9091}{rm}{61}\RG{148.170}{-0.004}{10.9091}{rm}{48}\RG{153.649}{-0.004}{10.9091}{mi}{58}%
}
\]
由归纳假设，存在\allowbreak{}\hskip.1em \RB{18.515}{8.186}{2.621}{%
\RG{0.000}{-0.005}{10.9091}{mi}{81}\RG{11.655}{-0.005}{10.9091}{sy}{50}%
}\hskip.1em\allowbreak{} GL\allowbreak{}\hskip.1em \RB{34.398}{8.677}{3.232}{%
\RG{0.000}{-1.642}{7.9701}{mi}{110}\RG{5.138}{-1.642}{7.9701}{sy}{0}\RG{11.724}{-1.642}{7.9701}{rm}{49}\RG{16.457}{-0.005}{10.9091}{rm}{40}\RG{20.699}{-0.005}{10.9091}{mi}{75}\RG{30.745}{-0.005}{10.9091}{rm}{41}%
}\hskip.1em\allowbreak{} ，使\allowbreak{}\hskip.1em \RB{37.331}{9.754}{2.621}{%
\RG{0.000}{-0.005}{10.9091}{mi}{81}\RG{8.624}{3.954}{7.9701}{sy}{0}\RG{15.211}{3.954}{7.9701}{rm}{49}\RG{19.943}{-0.005}{10.9091}{mi}{66}\RG{28.758}{-0.005}{10.9091}{mi}{81}%
}\hskip.1em\allowbreak{} 的主对角元全为\allowbreak{}\hskip.1em \RB{5.518}{7.761}{0.745}{%
\RG{0.000}{-0.005}{10.9091}{rm}{48}%
}\hskip.1em\allowbreak{} 。令\allowbreak{}\hskip.1em \RB{19.763}{7.946}{0.745}{%
\RG{0.000}{-0.005}{10.9091}{mi}{82}\RG{11.397}{-0.005}{10.9091}{rm}{61}%
}\hskip.1em\allowbreak{} diag\allowbreak{}\hskip.1em \RB{26.822}{8.677}{3.232}{%
\RG{0.000}{-0.005}{10.9091}{rm}{40}\RG{4.233}{-0.005}{10.9091}{rm}{49}\RG{9.697}{-0.005}{10.9091}{mi}{59}\RG{14.530}{-0.005}{10.9091}{mi}{81}\RG{23.169}{-0.005}{10.9091}{rm}{41}%
}\hskip.1em\allowbreak{} ，则

\[
\RB{180.517}{19.593}{14.124}{%
\RG{84.790}{18.656}{10.9092}{ex}{32}\RG{167.467}{18.656}{10.9092}{ex}{33}\RG{104.044}{7.281}{10.9091}{rm}{48}\RG{140.951}{7.281}{10.9091}{mi}{11}\RG{147.955}{7.281}{10.9091}{mi}{81}\RG{25.372}{4.506}{7.9701}{sy}{0}\RG{31.958}{4.506}{7.9701}{rm}{49}\RG{0.000}{0.001}{10.9091}{rm}{40}\RG{4.243}{0.001}{10.9091}{mi}{80}\RG{12.752}{0.001}{10.9091}{mi}{82}\RG{21.129}{0.001}{10.9091}{rm}{41}\RG{36.690}{0.001}{10.9091}{mi}{65}\RG{44.872}{0.001}{10.9091}{rm}{40}\RG{49.115}{0.001}{10.9091}{mi}{80}\RG{57.625}{0.001}{10.9091}{mi}{82}\RG{66.002}{0.001}{10.9091}{rm}{41}\RG{73.256}{0.001}{10.9091}{rm}{61}\RG{177.922}{0.001}{10.9091}{mi}{58}\RG{102.051}{-3.935}{7.9701}{sy}{0}\RG{108.637}{-3.935}{7.9701}{rm}{49}\RG{138.702}{-3.935}{7.9701}{sy}{0}\RG{145.289}{-3.935}{7.9701}{rm}{49}\RG{93.426}{-7.894}{10.9091}{mi}{81}\RG{113.369}{-7.894}{10.9091}{mi}{12}\RG{130.071}{-7.894}{10.9091}{mi}{81}\RG{150.021}{-7.894}{10.9091}{mi}{66}\RG{158.836}{-7.894}{10.9091}{mi}{81}%
}
\]
右端的主对角元全为\allowbreak{}\hskip.1em \RB{5.518}{7.767}{0.739}{%
\RG{0.000}{0.001}{10.9091}{rm}{48}%
}\hskip.1em\allowbreak{} ，归纳完成。\allowbreak{}\hskip.1em \RB{8.366}{7.996}{0.500}{%
\RG{0.000}{0.001}{10.9091}{am}{3}%
}\hskip.1em\allowbreak{}

小结：先选基使第一个对角元为零，再利用迹为零对剩余的主子块归纳。

题目来源：谢启鸿《复旦大学高等代数习题集（2023版）》，第83页。

% Original page 2.
\section*{2020S01}
设数域\allowbreak{}\hskip.1em \RB{10.198}{7.952}{0.500}{%
\RG{0.000}{0.001}{10.9091}{mi}{75}%
}\hskip.1em\allowbreak{} 上的二元多项式\allowbreak{}\hskip.1em \RB{31.234}{8.683}{3.226}{%
\RG{0.000}{0.001}{10.9091}{mi}{102}\RG{6.516}{0.001}{10.9091}{rm}{40}\RG{10.758}{0.001}{10.9091}{mi}{120}\RG{16.987}{0.001}{10.9091}{mi}{59}\RG{21.820}{0.001}{10.9091}{mi}{121}\RG{27.581}{0.001}{10.9091}{rm}{41}%
}\hskip.1em\allowbreak{} 关于\allowbreak{}\hskip.1em \RB{6.249}{5.323}{0.619}{%
\RG{0.000}{0.001}{10.9091}{mi}{120}%
}\hskip.1em\allowbreak{} 的次数不超过\allowbreak{}\hskip.1em \RB{6.729}{5.323}{0.619}{%
\RG{0.000}{0.001}{10.9091}{mi}{110}%
}\hskip.1em\allowbreak{} ，关于\allowbreak{}\hskip.1em \RB{5.846}{5.323}{2.735}{%
\RG{0.000}{0.001}{10.9091}{mi}{121}%
}\hskip.1em\allowbreak{} 的次数不超过\allowbreak{}\hskip.1em \RB{9.751}{5.323}{0.619}{%
\RG{0.000}{0.001}{10.9091}{mi}{109}%
}\hskip.1em\allowbreak{} 。设\allowbreak{}\hskip.1em \RB{10.198}{7.952}{0.500}{%
\RG{0.000}{0.001}{10.9091}{mi}{75}%
}\hskip.1em\allowbreak{} 中有两组数\allowbreak{}\hskip.1em \RB{61.233}{5.318}{2.609}{%
\RG{0.000}{-0.004}{10.9091}{mi}{97}\RG{5.766}{-1.640}{7.9701}{rm}{48}\RG{10.499}{-0.004}{10.9091}{mi}{59}\RG{15.332}{-0.004}{10.9091}{mi}{97}\RG{21.114}{-1.640}{7.9701}{rm}{49}\RG{25.846}{-0.004}{10.9091}{mi}{59}\RG{30.679}{-0.004}{10.9091}{mi}{58}\RG{35.523}{-0.004}{10.9091}{mi}{58}\RG{40.355}{-0.004}{10.9091}{mi}{58}\RG{45.199}{-0.004}{10.9091}{mi}{59}\RG{50.043}{-0.004}{10.9091}{mi}{97}\RG{55.855}{-1.640}{7.9701}{mi}{110}%
}\hskip.1em\allowbreak{} 和\allowbreak{}\hskip.1em \RB{60.330}{8.067}{2.609}{%
\RG{0.000}{-0.004}{10.9091}{mi}{98}\RG{4.682}{-1.640}{7.9701}{rm}{48}\RG{9.415}{-0.004}{10.9091}{mi}{59}\RG{14.248}{-0.004}{10.9091}{mi}{98}\RG{18.945}{-1.640}{7.9701}{rm}{49}\RG{23.677}{-0.004}{10.9091}{mi}{59}\RG{28.510}{-0.004}{10.9091}{mi}{58}\RG{33.354}{-0.004}{10.9091}{mi}{58}\RG{38.197}{-0.004}{10.9091}{mi}{58}\RG{43.030}{-0.004}{10.9091}{mi}{59}\RG{47.874}{-0.004}{10.9091}{mi}{98}\RG{52.601}{-1.640}{7.9701}{mi}{109}%
}\hskip.1em\allowbreak{} ，各组内的数两两不同，且

\[
\RB{193.533}{8.684}{3.760}{%
\RG{0.000}{0.002}{10.9091}{mi}{102}\RG{6.516}{0.002}{10.9091}{rm}{40}\RG{10.758}{0.002}{10.9091}{mi}{97}\RG{19.906}{0.002}{10.9091}{mi}{59}\RG{24.739}{0.002}{10.9091}{mi}{98}\RG{33.818}{0.002}{10.9091}{rm}{41}\RG{41.073}{0.002}{10.9091}{rm}{61}\RG{52.582}{0.002}{10.9091}{rm}{48}\RG{79.855}{0.002}{10.9091}{rm}{40}\RG{84.088}{0.002}{10.9091}{rm}{48}\RG{92.607}{0.002}{10.9091}{sy}{20}\RG{104.122}{0.002}{10.9091}{mi}{105}\RG{110.911}{0.002}{10.9091}{sy}{20}\RG{122.426}{0.002}{10.9091}{mi}{110}\RG{128.972}{0.002}{10.9091}{mi}{59}\RG{137.455}{0.002}{10.9091}{rm}{48}\RG{145.940}{0.002}{10.9091}{sy}{20}\RG{157.455}{0.002}{10.9091}{mi}{106}\RG{165.602}{0.002}{10.9091}{sy}{20}\RG{177.118}{0.002}{10.9091}{mi}{109}\RG{186.697}{0.002}{10.9091}{rm}{41}\RG{190.939}{0.002}{10.9091}{mi}{58}\RG{16.524}{-1.634}{7.9701}{mi}{105}\RG{29.436}{-1.634}{7.9701}{mi}{106}%
}
\]
证明：\allowbreak{}\hskip.1em \RB{31.234}{8.680}{3.229}{%
\RG{0.000}{-0.002}{10.9091}{mi}{102}\RG{6.515}{-0.002}{10.9091}{rm}{40}\RG{10.758}{-0.002}{10.9091}{mi}{120}\RG{16.987}{-0.002}{10.9091}{mi}{59}\RG{21.820}{-0.002}{10.9091}{mi}{121}\RG{27.581}{-0.002}{10.9091}{rm}{41}%
}\hskip.1em\allowbreak{} 是零多项式。

证明：按\allowbreak{}\hskip.1em \RB{6.249}{5.319}{0.623}{%
\RG{0.000}{-0.003}{10.9091}{mi}{120}%
}\hskip.1em\allowbreak{} 的幂次整理，写成

\[
\RB{255.896}{17.649}{14.023}{%
\RG{51.678}{13.634}{7.9701}{mi}{110}\RG{46.369}{10.361}{10.9092}{ex}{88}\RG{94.246}{4.502}{7.9701}{mi}{107}\RG{0.000}{-0.003}{10.9091}{mi}{102}\RG{6.515}{-0.003}{10.9091}{rm}{40}\RG{10.757}{-0.003}{10.9091}{mi}{120}\RG{16.986}{-0.003}{10.9091}{mi}{59}\RG{21.819}{-0.003}{10.9091}{mi}{121}\RG{27.581}{-0.003}{10.9091}{rm}{41}\RG{34.836}{-0.003}{10.9091}{rm}{61}\RG{63.945}{-0.003}{10.9091}{mi}{99}\RG{73.786}{-0.003}{10.9091}{rm}{40}\RG{78.028}{-0.003}{10.9091}{mi}{121}\RG{83.768}{-0.003}{10.9091}{rm}{41}\RG{88.011}{-0.003}{10.9091}{mi}{120}\RG{99.365}{-0.003}{10.9091}{mi}{59}\RG{126.016}{-0.003}{10.9091}{mi}{99}\RG{135.873}{-0.003}{10.9091}{rm}{40}\RG{140.115}{-0.003}{10.9091}{mi}{121}\RG{145.855}{-0.003}{10.9091}{rm}{41}\RG{153.128}{-0.003}{10.9091}{sy}{50}\RG{163.431}{-0.003}{10.9091}{mi}{75}\RG{173.476}{-0.003}{10.9091}{rm}{91}\RG{176.507}{-0.003}{10.9091}{mi}{121}\RG{182.248}{-0.003}{10.9091}{rm}{93}\RG{185.278}{-0.003}{10.9091}{mi}{59}\RT{201.035}{-0.003}{10.9091}{d}\RT{207.100}{-0.003}{10.9091}{e}\RT{211.944}{-0.003}{10.9091}{g}\RG{219.217}{-0.003}{10.9091}{mi}{99}\RG{232.088}{-0.003}{10.9091}{sy}{20}\RG{243.603}{-0.003}{10.9091}{mi}{109}\RG{253.181}{-0.003}{10.9091}{mi}{59}\RG{68.666}{-1.780}{7.9701}{mi}{107}\RG{130.753}{-1.780}{7.9701}{mi}{107}\RG{223.938}{-1.780}{7.9701}{mi}{107}\RG{46.526}{-13.356}{7.9701}{mi}{107}\RG{51.148}{-13.356}{7.9701}{rm}{61}\RG{57.732}{-13.356}{7.9701}{rm}{48}%
}
\]
其中允许某些\allowbreak{}\hskip.1em \RB{9.469}{5.323}{2.356}{%
\RG{0.000}{0.001}{10.9091}{mi}{99}\RG{4.721}{-1.776}{7.9701}{mi}{107}%
}\hskip.1em\allowbreak{} 为零多项式。

固定\allowbreak{}\hskip.1em \RB{78.642}{8.679}{3.230}{%
\RG{0.000}{-0.003}{10.9091}{mi}{106}\RG{8.147}{-0.003}{10.9091}{sy}{50}\RG{18.434}{-0.003}{10.9091}{sy}{102}\RG{23.905}{-0.003}{10.9091}{rm}{48}\RG{29.359}{-0.003}{10.9091}{mi}{59}\RG{34.208}{-0.003}{10.9091}{rm}{49}\RG{39.662}{-0.003}{10.9091}{mi}{59}\RG{44.495}{-0.003}{10.9091}{mi}{58}\RG{49.339}{-0.003}{10.9091}{mi}{58}\RG{54.182}{-0.003}{10.9091}{mi}{58}\RG{59.015}{-0.003}{10.9091}{mi}{59}\RG{63.859}{-0.003}{10.9091}{mi}{109}\RG{73.484}{-0.003}{10.9091}{sy}{103}%
}\hskip.1em\allowbreak{} 。一元多项式

\[
\RB{105.641}{17.647}{14.024}{%
\RG{55.002}{13.632}{7.9701}{mi}{110}\RG{49.692}{10.360}{10.9092}{ex}{88}\RG{100.893}{4.500}{7.9701}{mi}{107}\RG{0.000}{-0.004}{10.9091}{mi}{102}\RG{6.515}{-0.004}{10.9091}{rm}{40}\RG{10.757}{-0.004}{10.9091}{mi}{120}\RG{16.986}{-0.004}{10.9091}{mi}{59}\RG{21.819}{-0.004}{10.9091}{mi}{98}\RG{30.904}{-0.004}{10.9091}{rm}{41}\RG{38.159}{-0.004}{10.9091}{rm}{61}\RG{67.269}{-0.004}{10.9091}{mi}{99}\RG{77.110}{-0.004}{10.9091}{rm}{40}\RG{81.352}{-0.004}{10.9091}{mi}{98}\RG{90.416}{-0.004}{10.9091}{rm}{41}\RG{94.659}{-0.004}{10.9091}{mi}{120}\RG{26.522}{-1.640}{7.9701}{mi}{106}\RG{86.034}{-1.640}{7.9701}{mi}{106}\RG{71.990}{-1.781}{7.9701}{mi}{107}\RG{49.850}{-13.357}{7.9701}{mi}{107}\RG{54.471}{-13.357}{7.9701}{rm}{61}\RG{61.055}{-13.357}{7.9701}{rm}{48}%
}
\]
的次数不超过\allowbreak{}\hskip.1em \RB{6.729}{5.320}{0.622}{%
\RG{0.000}{-0.002}{10.9091}{mi}{110}%
}\hskip.1em\allowbreak{} ，且有\allowbreak{}\hskip.1em \RB{25.378}{7.764}{1.407}{%
\RG{0.000}{-0.002}{10.9091}{mi}{110}\RG{9.191}{-0.002}{10.9091}{rm}{43}\RG{20.308}{-0.002}{10.9091}{rm}{49}%
}\hskip.1em\allowbreak{} 个不同的根\allowbreak{}\hskip.1em \RB{61.232}{5.320}{2.607}{%
\RG{0.000}{-0.002}{10.9091}{mi}{97}\RG{5.766}{-1.639}{7.9701}{rm}{48}\RG{10.498}{-0.002}{10.9091}{mi}{59}\RG{15.331}{-0.002}{10.9091}{mi}{97}\RG{21.113}{-1.639}{7.9701}{rm}{49}\RG{25.846}{-0.002}{10.9091}{mi}{59}\RG{30.679}{-0.002}{10.9091}{mi}{58}\RG{35.523}{-0.002}{10.9091}{mi}{58}\RG{40.355}{-0.002}{10.9091}{mi}{58}\RG{45.199}{-0.002}{10.9091}{mi}{59}\RG{50.043}{-0.002}{10.9091}{mi}{97}\RG{55.854}{-1.639}{7.9701}{mi}{110}%
}\hskip.1em\allowbreak{} 。非零一元多项式的不同根数不超过其次数，因此

\[
\RB{57.742}{8.680}{3.765}{%
\RG{0.000}{-0.002}{10.9091}{mi}{102}\RG{6.515}{-0.002}{10.9091}{rm}{40}\RG{10.758}{-0.002}{10.9091}{mi}{120}\RG{16.987}{-0.002}{10.9091}{mi}{59}\RG{21.820}{-0.002}{10.9091}{mi}{98}\RG{30.905}{-0.002}{10.9091}{rm}{41}\RG{38.178}{-0.002}{10.9091}{sy}{17}\RG{49.693}{-0.002}{10.9091}{rm}{48}\RG{55.147}{-0.002}{10.9091}{mi}{58}\RG{26.523}{-1.639}{7.9701}{mi}{106}%
}
\]
比较\allowbreak{}\hskip.1em \RB{6.249}{5.319}{0.623}{%
\RG{0.000}{-0.003}{10.9091}{mi}{120}%
}\hskip.1em\allowbreak{} 的各次幂的系数，得到

\[
\RB{185.127}{8.685}{3.759}{%
\RG{0.000}{0.003}{10.9091}{mi}{99}\RG{9.841}{0.003}{10.9091}{rm}{40}\RG{14.084}{0.003}{10.9091}{mi}{98}\RG{23.148}{0.003}{10.9091}{rm}{41}\RG{30.403}{0.003}{10.9091}{rm}{61}\RG{41.912}{0.003}{10.9091}{rm}{48}\RG{69.185}{0.003}{10.9091}{rm}{40}\RG{73.418}{0.003}{10.9091}{rm}{48}\RG{81.936}{0.003}{10.9091}{sy}{20}\RG{93.452}{0.003}{10.9091}{mi}{107}\RG{102.505}{0.003}{10.9091}{sy}{20}\RG{114.020}{0.003}{10.9091}{mi}{110}\RG{120.565}{0.003}{10.9091}{mi}{59}\RG{129.049}{0.003}{10.9091}{rm}{48}\RG{137.534}{0.003}{10.9091}{sy}{20}\RG{149.049}{0.003}{10.9091}{mi}{106}\RG{157.196}{0.003}{10.9091}{sy}{20}\RG{168.711}{0.003}{10.9091}{mi}{109}\RG{178.289}{0.003}{10.9091}{rm}{41}\RG{182.533}{0.003}{10.9091}{mi}{58}\RG{18.765}{-1.633}{7.9701}{mi}{106}\RG{4.721}{-1.774}{7.9701}{mi}{107}%
}
\]
再固定\allowbreak{}\hskip.1em \RB{76.518}{8.682}{3.227}{%
\RG{0.000}{0.000}{10.9091}{mi}{107}\RG{9.053}{0.000}{10.9091}{sy}{50}\RG{19.340}{0.000}{10.9091}{sy}{102}\RG{24.811}{0.000}{10.9091}{rm}{48}\RG{30.265}{0.000}{10.9091}{mi}{59}\RG{35.114}{0.000}{10.9091}{rm}{49}\RG{40.568}{0.000}{10.9091}{mi}{59}\RG{45.401}{0.000}{10.9091}{mi}{58}\RG{50.245}{0.000}{10.9091}{mi}{58}\RG{55.088}{0.000}{10.9091}{mi}{58}\RG{59.921}{0.000}{10.9091}{mi}{59}\RG{64.765}{0.000}{10.9091}{mi}{110}\RG{71.360}{0.000}{10.9091}{sy}{103}%
}\hskip.1em\allowbreak{} 。多项式\allowbreak{}\hskip.1em \RB{23.476}{8.682}{3.227}{%
\RG{0.000}{0.000}{10.9091}{mi}{99}\RG{4.721}{-1.777}{7.9701}{mi}{107}\RG{9.841}{0.000}{10.9091}{rm}{40}\RG{14.083}{0.000}{10.9091}{mi}{121}\RG{19.823}{0.000}{10.9091}{rm}{41}%
}\hskip.1em\allowbreak{} 的次数不超过\allowbreak{}\hskip.1em \RB{9.751}{5.322}{0.620}{%
\RG{0.000}{0.000}{10.9091}{mi}{109}%
}\hskip.1em\allowbreak{} ，而\allowbreak{}\hskip.1em \RB{60.330}{8.071}{2.605}{%
\RG{0.000}{0.000}{10.9091}{mi}{98}\RG{4.682}{-1.636}{7.9701}{rm}{48}\RG{9.414}{0.000}{10.9091}{mi}{59}\RG{14.247}{0.000}{10.9091}{mi}{98}\RG{18.944}{-1.636}{7.9701}{rm}{49}\RG{23.677}{0.000}{10.9091}{mi}{59}\RG{28.510}{0.000}{10.9091}{mi}{58}\RG{33.354}{0.000}{10.9091}{mi}{58}\RG{38.186}{0.000}{10.9091}{mi}{58}\RG{43.030}{0.000}{10.9091}{mi}{59}\RG{47.874}{0.000}{10.9091}{mi}{98}\RG{52.601}{-1.636}{7.9701}{mi}{109}%
}\hskip.1em\allowbreak{} 是它的\allowbreak{}\hskip.1em \RB{27.280}{7.766}{1.405}{%
\RG{0.000}{0.000}{10.9091}{mi}{109}\RG{11.660}{0.000}{10.9091}{rm}{43}\RG{22.209}{0.000}{10.9091}{rm}{49}%
}\hskip.1em\allowbreak{} 个不同的根。因此\allowbreak{}\hskip.1em \RB{44.130}{8.677}{3.232}{%
\RG{0.000}{-0.005}{10.9091}{mi}{99}\RG{4.721}{-1.782}{7.9701}{mi}{107}\RG{9.841}{-0.005}{10.9091}{rm}{40}\RG{14.083}{-0.005}{10.9091}{mi}{121}\RG{19.823}{-0.005}{10.9091}{rm}{41}\RG{27.096}{-0.005}{10.9091}{sy}{17}\RG{38.611}{-0.005}{10.9091}{rm}{48}%
}\hskip.1em\allowbreak{} 。于是所有系数多项式都为零，即

\[
\RB{54.418}{8.683}{3.226}{%
\RG{0.000}{0.001}{10.9091}{mi}{102}\RG{6.515}{0.001}{10.9091}{rm}{40}\RG{10.758}{0.001}{10.9091}{mi}{120}\RG{16.987}{0.001}{10.9091}{mi}{59}\RG{21.820}{0.001}{10.9091}{mi}{121}\RG{27.581}{0.001}{10.9091}{rm}{41}\RG{34.854}{0.001}{10.9091}{sy}{17}\RG{46.369}{0.001}{10.9091}{rm}{48}\RG{51.823}{0.001}{10.9091}{mi}{58}%
}
\]
上述根数结论可由因式定理直接得到：若非零多项式有\allowbreak{}\hskip.1em \RB{5.256}{5.319}{0.623}{%
\RG{0.000}{-0.003}{10.9091}{mi}{114}%
}\hskip.1em\allowbreak{} 个不同的根，则对应的\allowbreak{}\hskip.1em \RB{5.256}{5.319}{0.623}{%
\RG{0.000}{-0.003}{10.9091}{mi}{114}%
}\hskip.1em\allowbreak{} 个一次因子的乘积整除它，故其次数至少为\allowbreak{}\hskip.1em \RB{5.256}{5.324}{0.618}{%
\RG{0.000}{0.002}{10.9091}{mi}{114}%
}\hskip.1em\allowbreak{} 。\allowbreak{}\hskip.1em \RB{8.366}{7.997}{0.500}{%
\RG{0.000}{0.002}{10.9091}{am}{3}%
}\hskip.1em\allowbreak{}

小结：依次固定一个变量，两次使用一元多项式的根数上界；关键是分别控制两个变量的次数。

题目来源：谢启鸿《复旦大学高等代数习题集（2023版）》，第93页。

% Original page 3.
\section*{2020S02}
用特征值法证明下列结论（分别对应高代白皮书例5.76、例3.82和例3.84）。

(1)设\allowbreak{}\hskip.1em \RB{81.277}{8.679}{3.230}{%
\RG{0.000}{-0.003}{10.9091}{mi}{102}\RG{6.516}{-0.003}{10.9091}{rm}{40}\RG{10.758}{-0.003}{10.9091}{mi}{120}\RG{16.993}{-0.003}{10.9091}{rm}{41}\RG{21.235}{-0.003}{10.9091}{mi}{59}\RG{26.068}{-0.003}{10.9091}{mi}{103}\RG{31.679}{-0.003}{10.9091}{rm}{40}\RG{35.921}{-0.003}{10.9091}{mi}{120}\RG{42.156}{-0.003}{10.9091}{rm}{41}\RG{49.429}{-0.003}{10.9091}{sy}{50}\RG{59.731}{-0.003}{10.9091}{mi}{75}\RG{69.778}{-0.003}{10.9091}{rm}{91}\RG{72.808}{-0.003}{10.9091}{mi}{120}\RG{79.043}{-0.003}{10.9091}{rm}{93}%
}\hskip.1em\allowbreak{} 互素，\allowbreak{}\hskip.1em \RB{55.676}{8.679}{3.230}{%
\RG{0.000}{-0.003}{10.9091}{mi}{65}\RG{11.212}{-0.003}{10.9091}{sy}{50}\RG{21.515}{-0.003}{10.9091}{mi}{77}\RG{32.099}{-1.639}{7.9701}{mi}{110}\RG{37.735}{-0.003}{10.9091}{rm}{40}\RG{41.978}{-0.003}{10.9091}{mi}{75}\RG{52.023}{-0.003}{10.9091}{rm}{41}%
}\hskip.1em\allowbreak{} 且\allowbreak{}\hskip.1em \RB{43.222}{8.679}{3.230}{%
\RG{0.000}{-0.003}{10.9091}{mi}{102}\RG{6.515}{-0.003}{10.9091}{rm}{40}\RG{10.758}{-0.003}{10.9091}{mi}{65}\RG{18.940}{-0.003}{10.9091}{rm}{41}\RG{26.195}{-0.003}{10.9091}{rm}{61}\RG{37.704}{-0.003}{10.9091}{rm}{48}%
}\hskip.1em\allowbreak{} 。证明：\allowbreak{}\hskip.1em \RB{21.672}{8.679}{3.230}{%
\RG{0.000}{-0.003}{10.9091}{mi}{103}\RG{5.595}{-0.003}{10.9091}{rm}{40}\RG{9.837}{-0.003}{10.9091}{mi}{65}\RG{18.019}{-0.003}{10.9091}{rm}{41}%
}\hskip.1em\allowbreak{} 非奇异。(2)设\allowbreak{}\hskip.1em \RB{8.366}{8.315}{0.500}{%
\RG{0.000}{0.004}{10.9091}{mi}{65}%
}\hskip.1em\allowbreak{} 为\allowbreak{}\hskip.1em \RB{6.729}{5.326}{0.616}{%
\RG{0.000}{0.004}{10.9091}{mi}{110}%
}\hskip.1em\allowbreak{} 阶实对称矩阵，\allowbreak{}\hskip.1em \RB{7.536}{8.195}{0.736}{%
\RG{0.000}{0.004}{10.9091}{mi}{83}%
}\hskip.1em\allowbreak{} 为\allowbreak{}\hskip.1em \RB{6.729}{5.326}{0.616}{%
\RG{0.000}{0.004}{10.9091}{mi}{110}%
}\hskip.1em\allowbreak{} 阶实反对称矩阵。证明：\allowbreak{}\hskip.1em \RB{31.302}{8.195}{2.213}{%
\RG{0.000}{0.004}{10.9091}{mi}{73}\RG{4.795}{-1.633}{7.9701}{mi}{110}\RG{12.856}{0.004}{10.9091}{sy}{6}\RG{23.765}{0.004}{10.9091}{mi}{83}%
}\hskip.1em\allowbreak{} 与\allowbreak{}\hskip.1em \RB{35.883}{8.315}{2.213}{%
\RG{0.000}{0.004}{10.9091}{mi}{73}\RG{4.796}{-1.633}{7.9701}{mi}{110}\RG{12.856}{0.004}{10.9091}{sy}{6}\RG{23.765}{0.004}{10.9091}{mi}{105}\RG{27.518}{0.004}{10.9091}{mi}{65}%
}\hskip.1em\allowbreak{} 都非奇异。

证明：(1)在\allowbreak{}\hskip.1em \RB{7.973}{8.182}{0.705}{%
\RG{0.000}{0.002}{10.9091}{bb}{67}%
}\hskip.1em\allowbreak{} 上考虑\allowbreak{}\hskip.1em \RB{8.366}{8.313}{0.500}{%
\RG{0.000}{0.002}{10.9091}{mi}{65}%
}\hskip.1em\allowbreak{} ，设其特征值（计重数）为\allowbreak{}\hskip.1em \RB{47.080}{8.073}{2.603}{%
\RG{0.000}{0.002}{10.9091}{mi}{21}\RG{6.363}{-1.634}{7.9701}{rm}{49}\RG{11.096}{0.002}{10.9091}{mi}{59}\RG{15.929}{0.002}{10.9091}{mi}{58}\RG{20.773}{0.002}{10.9091}{mi}{58}\RG{25.605}{0.002}{10.9091}{mi}{58}\RG{30.449}{0.002}{10.9091}{mi}{59}\RG{35.293}{0.002}{10.9091}{mi}{21}\RG{41.702}{-1.634}{7.9701}{mi}{110}%
}\hskip.1em\allowbreak{} 。若\allowbreak{}\hskip.1em \RB{44.758}{8.313}{3.760}{%
\RG{0.000}{0.002}{10.9091}{mi}{65}\RG{8.182}{0.002}{10.9091}{mi}{118}\RG{16.891}{0.002}{10.9091}{rm}{61}\RG{28.406}{0.002}{10.9091}{mi}{21}\RG{34.770}{-1.634}{7.9701}{mi}{106}\RG{39.152}{0.002}{10.9091}{mi}{118}%
}\hskip.1em\allowbreak{} 且\allowbreak{}\hskip.1em \RB{25.727}{8.313}{2.854}{%
\RG{0.000}{0.002}{10.9091}{mi}{118}\RG{8.710}{0.002}{10.9091}{sy}{54}\RG{8.710}{0.002}{10.9091}{rm}{61}\RG{20.208}{0.002}{10.9091}{rm}{48}%
}\hskip.1em\allowbreak{} ，则

\[
\RB{97.529}{8.680}{3.764}{%
\RG{0.000}{-0.002}{10.9091}{rm}{48}\RG{8.476}{-0.002}{10.9091}{rm}{61}\RG{20.000}{-0.002}{10.9091}{mi}{102}\RG{26.515}{-0.002}{10.9091}{rm}{40}\RG{30.757}{-0.002}{10.9091}{mi}{65}\RG{38.939}{-0.002}{10.9091}{rm}{41}\RG{43.182}{-0.002}{10.9091}{mi}{118}\RG{51.891}{-0.002}{10.9091}{rm}{61}\RG{63.407}{-0.002}{10.9091}{mi}{102}\RG{69.923}{-0.002}{10.9091}{rm}{40}\RG{74.165}{-0.002}{10.9091}{mi}{21}\RG{84.911}{-0.002}{10.9091}{rm}{41}\RG{89.153}{-0.002}{10.9091}{mi}{118}\RG{94.815}{-0.002}{10.9091}{mi}{59}\RG{80.529}{-1.638}{7.9701}{mi}{106}%
}
\]
故\allowbreak{}\hskip.1em \RB{45.785}{8.686}{3.758}{%
\RG{0.000}{0.004}{10.9091}{mi}{102}\RG{6.515}{0.004}{10.9091}{rm}{40}\RG{10.757}{0.004}{10.9091}{mi}{21}\RG{17.121}{-1.632}{7.9701}{mi}{106}\RG{21.503}{0.004}{10.9091}{rm}{41}\RG{28.758}{0.004}{10.9091}{rm}{61}\RG{40.267}{0.004}{10.9091}{rm}{48}%
}\hskip.1em\allowbreak{} 。由\allowbreak{}\hskip.1em \RB{16.406}{8.195}{2.732}{%
\RG{0.000}{0.004}{10.9091}{mi}{102}\RG{5.902}{0.004}{10.9091}{mi}{59}\RG{10.735}{0.004}{10.9091}{mi}{103}%
}\hskip.1em\allowbreak{} 互素，存在\allowbreak{}\hskip.1em \RB{50.477}{8.686}{3.223}{%
\RG{0.000}{0.004}{10.9091}{mi}{112}\RG{5.487}{0.004}{10.9091}{mi}{59}\RG{10.320}{0.004}{10.9091}{mi}{113}\RG{18.628}{0.004}{10.9091}{sy}{50}\RG{28.931}{0.004}{10.9091}{mi}{75}\RG{38.977}{0.004}{10.9091}{rm}{91}\RG{42.007}{0.004}{10.9091}{mi}{120}\RG{48.242}{0.004}{10.9091}{rm}{93}%
}\hskip.1em\allowbreak{} ，使

\[
\RB{117.668}{8.682}{3.227}{%
\RG{0.000}{0.000}{10.9091}{mi}{112}\RG{5.488}{0.000}{10.9091}{rm}{40}\RG{9.731}{0.000}{10.9091}{mi}{120}\RG{15.966}{0.000}{10.9091}{rm}{41}\RG{20.208}{0.000}{10.9091}{mi}{102}\RG{26.723}{0.000}{10.9091}{rm}{40}\RG{30.966}{0.000}{10.9091}{mi}{120}\RG{37.201}{0.000}{10.9091}{rm}{41}\RG{43.856}{0.000}{10.9091}{rm}{43}\RG{54.776}{0.000}{10.9091}{mi}{113}\RG{60.038}{0.000}{10.9091}{rm}{40}\RG{64.281}{0.000}{10.9091}{mi}{120}\RG{70.516}{0.000}{10.9091}{rm}{41}\RG{74.758}{0.000}{10.9091}{mi}{103}\RG{80.353}{0.000}{10.9091}{rm}{40}\RG{84.596}{0.000}{10.9091}{mi}{120}\RG{90.830}{0.000}{10.9091}{rm}{41}\RG{98.084}{0.000}{10.9091}{rm}{61}\RG{109.594}{0.000}{10.9091}{rm}{49}\RG{115.073}{0.000}{10.9091}{mi}{58}%
}
\]
代入\allowbreak{}\hskip.1em \RB{30.936}{8.067}{3.766}{%
\RG{0.000}{-0.004}{10.9091}{mi}{120}\RG{9.265}{-0.004}{10.9091}{rm}{61}\RG{20.780}{-0.004}{10.9091}{mi}{21}\RG{27.144}{-1.640}{7.9701}{mi}{106}%
}\hskip.1em\allowbreak{} ，得\allowbreak{}\hskip.1em \RB{68.910}{8.678}{3.766}{%
\RG{0.000}{-0.004}{10.9091}{mi}{113}\RG{5.261}{-0.004}{10.9091}{rm}{40}\RG{9.504}{-0.004}{10.9091}{mi}{21}\RG{15.867}{-1.640}{7.9701}{mi}{106}\RG{20.249}{-0.004}{10.9091}{rm}{41}\RG{24.492}{-0.004}{10.9091}{mi}{103}\RG{30.087}{-0.004}{10.9091}{rm}{40}\RG{34.329}{-0.004}{10.9091}{mi}{21}\RG{40.693}{-1.640}{7.9701}{mi}{106}\RG{45.075}{-0.004}{10.9091}{rm}{41}\RG{52.330}{-0.004}{10.9091}{rm}{61}\RG{63.839}{-0.004}{10.9091}{rm}{49}%
}\hskip.1em\allowbreak{} ，所以\allowbreak{}\hskip.1em \RB{44.873}{8.678}{3.766}{%
\RG{0.000}{-0.004}{10.9091}{mi}{103}\RG{5.595}{-0.004}{10.9091}{rm}{40}\RG{9.837}{-0.004}{10.9091}{mi}{21}\RG{16.201}{-1.640}{7.9701}{mi}{106}\RG{20.583}{-0.004}{10.9091}{rm}{41}\RG{27.856}{-0.004}{10.9091}{sy}{54}\RG{27.856}{-0.004}{10.9091}{rm}{61}\RG{39.355}{-0.004}{10.9091}{rm}{48}%
}\hskip.1em\allowbreak{} 。

将\allowbreak{}\hskip.1em \RB{8.366}{8.314}{0.500}{%
\RG{0.000}{0.003}{10.9091}{mi}{65}%
}\hskip.1em\allowbreak{} 在复数域上相似化为上三角矩阵\allowbreak{}\hskip.1em \RB{8.180}{7.889}{0.500}{%
\RG{0.000}{0.003}{10.9091}{mi}{84}%
}\hskip.1em\allowbreak{} ，其对角元为\allowbreak{}\hskip.1em \RB{47.080}{8.074}{2.602}{%
\RG{0.000}{0.003}{10.9091}{mi}{21}\RG{6.364}{-1.634}{7.9701}{rm}{49}\RG{11.096}{0.003}{10.9091}{mi}{59}\RG{15.929}{0.003}{10.9091}{mi}{58}\RG{20.773}{0.003}{10.9091}{mi}{58}\RG{25.616}{0.003}{10.9091}{mi}{58}\RG{30.449}{0.003}{10.9091}{mi}{59}\RG{35.293}{0.003}{10.9091}{mi}{21}\RG{41.702}{-1.634}{7.9701}{mi}{110}%
}\hskip.1em\allowbreak{} ，则\allowbreak{}\hskip.1em \RB{21.380}{8.685}{3.224}{%
\RG{0.000}{0.003}{10.9091}{mi}{103}\RG{5.595}{0.003}{10.9091}{rm}{40}\RG{9.837}{0.003}{10.9091}{mi}{84}\RG{17.727}{0.003}{10.9091}{rm}{41}%
}\hskip.1em\allowbreak{} 的对角元为\allowbreak{}\hskip.1em \RB{74.908}{8.680}{3.229}{%
\RG{0.000}{-0.002}{10.9091}{mi}{103}\RG{5.595}{-0.002}{10.9091}{rm}{40}\RG{9.837}{-0.002}{10.9091}{mi}{21}\RG{16.201}{-1.639}{7.9701}{rm}{49}\RG{20.933}{-0.002}{10.9091}{rm}{41}\RG{25.176}{-0.002}{10.9091}{mi}{59}\RG{30.009}{-0.002}{10.9091}{mi}{58}\RG{34.853}{-0.002}{10.9091}{mi}{58}\RG{39.685}{-0.002}{10.9091}{mi}{58}\RG{44.529}{-0.002}{10.9091}{mi}{59}\RG{49.373}{-0.002}{10.9091}{mi}{103}\RG{55.013}{-0.002}{10.9091}{rm}{40}\RG{59.255}{-0.002}{10.9091}{mi}{21}\RG{65.620}{-1.639}{7.9701}{mi}{110}\RG{71.256}{-0.002}{10.9091}{rm}{41}%
}\hskip.1em\allowbreak{} 。因此

\[
\RB{117.722}{17.657}{15.208}{%
\RG{58.561}{13.642}{7.9701}{mi}{110}\RG{54.160}{10.369}{10.9092}{ex}{89}\RT{0.000}{0.005}{10.9091}{d}\RT{6.066}{0.005}{10.9091}{e}\RT{10.909}{0.005}{10.9091}{t}\RG{16.971}{0.005}{10.9091}{mi}{103}\RG{22.565}{0.005}{10.9091}{rm}{40}\RG{26.808}{0.005}{10.9091}{mi}{65}\RG{34.990}{0.005}{10.9091}{rm}{41}\RG{42.245}{0.005}{10.9091}{rm}{61}\RG{70.301}{0.005}{10.9091}{mi}{103}\RG{75.896}{0.005}{10.9091}{rm}{40}\RG{80.139}{0.005}{10.9091}{mi}{21}\RG{90.884}{0.005}{10.9091}{rm}{41}\RG{98.157}{0.005}{10.9091}{sy}{54}\RG{98.157}{0.005}{10.9091}{rm}{61}\RG{109.655}{0.005}{10.9091}{rm}{48}\RG{115.127}{0.005}{10.9091}{mi}{58}\RG{86.502}{-1.631}{7.9701}{mi}{106}\RG{53.778}{-13.082}{7.9701}{mi}{106}\RG{57.662}{-13.082}{7.9701}{rm}{61}\RG{64.246}{-13.082}{7.9701}{rm}{49}%
}
\]
这里所用的上三角化可按维数归纳证明：先取一个特征向量作为首个基向量，再对表示矩阵右下角的低一阶方块作同样处理。

(2)对复向量\allowbreak{}\hskip.1em \RB{5.595}{5.319}{0.623}{%
\RG{0.000}{-0.003}{10.9091}{mi}{122}%
}\hskip.1em\allowbreak{} ，记\allowbreak{}\hskip.1em \RB{9.719}{8.162}{0.623}{%
\RG{0.000}{-0.003}{10.9091}{mi}{122}\RG{5.553}{3.956}{7.9701}{sy}{3}%
}\hskip.1em\allowbreak{} 为其共轭转置。若\allowbreak{}\hskip.1em \RB{39.372}{8.188}{0.743}{%
\RG{0.000}{-0.003}{10.9091}{mi}{83}\RG{7.309}{-0.003}{10.9091}{mi}{122}\RG{15.902}{-0.003}{10.9091}{rm}{61}\RG{27.417}{-0.003}{10.9091}{mi}{21}\RG{33.777}{-0.003}{10.9091}{mi}{122}%
}\hskip.1em\allowbreak{} 、\allowbreak{}\hskip.1em \RB{25.600}{8.308}{2.859}{%
\RG{0.000}{-0.003}{10.9091}{mi}{122}\RG{8.583}{-0.003}{10.9091}{sy}{54}\RG{8.583}{-0.003}{10.9091}{rm}{61}\RG{20.081}{-0.003}{10.9091}{rm}{48}%
}\hskip.1em\allowbreak{} ，由\allowbreak{}\hskip.1em \RB{42.617}{8.188}{0.743}{%
\RG{0.000}{-0.003}{10.9091}{mi}{83}\RG{7.318}{3.956}{7.9701}{sy}{3}\RG{15.080}{-0.003}{10.9091}{rm}{61}\RG{26.596}{-0.003}{10.9091}{sy}{0}\RG{35.080}{-0.003}{10.9091}{mi}{83}%
}\hskip.1em\allowbreak{} ，有

\[
\RB{114.363}{26.874}{0.740}{%
\RG{5.553}{22.668}{7.9701}{sy}{3}\RG{49.619}{22.668}{7.9701}{sy}{3}\RG{0.000}{18.163}{10.9091}{mi}{122}\RG{10.285}{18.163}{10.9091}{mi}{83}\RG{17.594}{18.163}{10.9091}{mi}{122}\RG{26.187}{18.163}{10.9091}{rm}{61}\RG{37.702}{18.163}{10.9091}{mi}{21}\RG{44.062}{18.163}{10.9091}{mi}{122}\RG{54.351}{18.163}{10.9091}{mi}{122}\RG{59.893}{18.163}{10.9091}{mi}{59}\RG{43.255}{4.504}{7.9701}{sy}{3}\RG{55.305}{4.504}{7.9701}{sy}{3}\RG{94.175}{4.504}{7.9701}{sy}{3}\RG{5.553}{3.151}{7.9701}{sy}{3}\RG{0.000}{-0.000}{10.9091}{mi}{122}\RG{10.285}{-0.000}{10.9091}{mi}{83}\RG{17.594}{-0.000}{10.9091}{mi}{122}\RG{26.187}{-0.000}{10.9091}{rm}{61}\RG{37.702}{-0.000}{10.9091}{mi}{122}\RG{47.987}{-0.000}{10.9091}{mi}{83}\RG{60.039}{-0.000}{10.9091}{mi}{122}\RG{68.622}{-0.000}{10.9091}{rm}{61}\RG{80.137}{-0.000}{10.9091}{sy}{0}\RG{88.622}{-0.000}{10.9091}{mi}{122}\RG{98.907}{-0.000}{10.9091}{mi}{83}\RG{106.216}{-0.000}{10.9091}{mi}{122}\RG{111.769}{-0.000}{10.9091}{mi}{58}\RR{0.000}{8.981}{23.156}{0.436}%
}
\]
因\allowbreak{}\hskip.1em \RB{9.806}{0.500}{16.954}{%
\RG{0.000}{-5.981}{10.9091}{sy}{112}%
}\hskip.1em\allowbreak{} 为\allowbreak{}\hskip.1em \RB{36.458}{8.161}{0.940}{%
\RG{0.000}{-0.004}{10.9091}{mi}{122}\RG{5.553}{3.955}{7.9701}{sy}{3}\RG{10.285}{-0.004}{10.9091}{mi}{122}\RG{19.143}{-0.004}{10.9091}{mi}{62}\RG{30.940}{-0.004}{10.9091}{rm}{48}%
}\hskip.1em\allowbreak{} ，所以\allowbreak{}\hskip.1em \RB{36.429}{8.067}{0.646}{%
\RG{0.000}{-0.004}{10.9091}{mi}{21}\RG{9.672}{-0.004}{10.9091}{rm}{61}\RG{21.466}{-0.004}{10.9091}{sy}{0}\RG{29.951}{-0.004}{10.9091}{mi}{21}\RR{0.000}{9.099}{6.364}{0.436}%
}\hskip.1em\allowbreak{} ，即\allowbreak{}\hskip.1em \RB{29.319}{8.067}{0.646}{%
\RG{0.000}{-0.004}{10.9091}{mi}{21}\RG{9.672}{-0.004}{10.9091}{rm}{61}\RG{21.466}{-0.004}{10.9091}{mi}{105}\RG{25.219}{-0.004}{10.9091}{mi}{116}%
}\hskip.1em\allowbreak{} ，其中\allowbreak{}\hskip.1em \RB{25.998}{7.969}{0.940}{%
\RG{0.000}{-0.004}{10.9091}{mi}{116}\RG{7.248}{-0.004}{10.9091}{sy}{50}\RG{17.829}{-0.004}{10.9091}{bb}{82}%
}\hskip.1em\allowbreak{} 。故\allowbreak{}\hskip.1em \RB{31.525}{8.187}{2.221}{%
\RG{0.000}{-0.004}{10.9091}{mi}{73}\RG{4.795}{-1.641}{7.9701}{mi}{110}\RG{12.967}{-0.004}{10.9091}{sy}{6}\RG{23.988}{-0.004}{10.9091}{mi}{83}%
}\hskip.1em\allowbreak{} 的特征值均形如\allowbreak{}\hskip.1em \RB{26.863}{7.762}{0.624}{%
\RG{0.000}{-0.004}{10.9091}{rm}{49}\RG{7.990}{-0.004}{10.9091}{sy}{6}\RG{19.010}{-0.004}{10.9091}{mi}{105}\RG{22.763}{-0.004}{10.9091}{mi}{116}%
}\hskip.1em\allowbreak{} ，其模为\allowbreak{}\hskip.1em \RB{47.524}{8.952}{1.404}{%
\RG{0.000}{0.001}{10.9091}{rm}{49}\RG{7.876}{0.001}{10.9091}{rm}{43}\RG{18.788}{0.001}{10.9091}{mi}{116}\RG{22.727}{3.152}{7.9701}{rm}{50}\RG{30.490}{0.001}{10.9091}{mi}{62}\RG{42.005}{0.001}{10.9091}{rm}{48}\RR{0.000}{9.417}{27.459}{0.436}%
}\hskip.1em\allowbreak{} 。

类似地，若\allowbreak{}\hskip.1em \RB{40.443}{8.308}{2.859}{%
\RG{0.000}{-0.003}{10.9091}{mi}{65}\RG{8.182}{-0.003}{10.9091}{mi}{122}\RG{16.765}{-0.003}{10.9091}{rm}{61}\RG{28.281}{-0.003}{10.9091}{mi}{22}\RG{34.849}{-0.003}{10.9091}{mi}{122}%
}\hskip.1em\allowbreak{} ，则由\allowbreak{}\hskip.1em \RB{35.825}{8.308}{0.503}{%
\RG{0.000}{-0.003}{10.9091}{mi}{65}\RG{8.182}{3.956}{7.9701}{sy}{3}\RG{15.944}{-0.003}{10.9091}{rm}{61}\RG{27.459}{-0.003}{10.9091}{mi}{65}%
}\hskip.1em\allowbreak{} 得

\[
\RB{65.290}{8.713}{2.602}{%
\RG{44.118}{4.507}{7.9701}{sy}{3}\RG{5.553}{3.154}{7.9701}{sy}{3}\RG{0.000}{0.003}{10.9091}{mi}{122}\RG{10.285}{0.003}{10.9091}{mi}{65}\RG{18.467}{0.003}{10.9091}{mi}{122}\RG{27.050}{0.003}{10.9091}{rm}{61}\RG{38.565}{0.003}{10.9091}{mi}{122}\RG{48.851}{0.003}{10.9091}{mi}{65}\RG{57.033}{0.003}{10.9091}{mi}{122}\RG{62.575}{0.003}{10.9091}{mi}{59}\RR{0.000}{8.985}{24.020}{0.436}%
}
\]
从而\allowbreak{}\hskip.1em \RB{28.075}{7.972}{2.857}{%
\RG{0.000}{-0.001}{10.9091}{mi}{22}\RG{9.603}{-0.001}{10.9091}{sy}{50}\RG{19.906}{-0.001}{10.9091}{bb}{82}%
}\hskip.1em\allowbreak{} 。于是\allowbreak{}\hskip.1em \RB{35.883}{8.310}{2.217}{%
\RG{0.000}{-0.001}{10.9091}{mi}{73}\RG{4.796}{-1.637}{7.9701}{mi}{110}\RG{12.856}{-0.001}{10.9091}{sy}{6}\RG{23.765}{-0.001}{10.9091}{mi}{105}\RG{27.518}{-0.001}{10.9091}{mi}{65}%
}\hskip.1em\allowbreak{} 的特征值均形如\allowbreak{}\hskip.1em \RB{29.281}{7.764}{2.857}{%
\RG{0.000}{-0.001}{10.9091}{rm}{49}\RG{7.879}{-0.001}{10.9091}{sy}{6}\RG{18.788}{-0.001}{10.9091}{mi}{105}\RG{22.541}{-0.001}{10.9091}{mi}{22}%
}\hskip.1em\allowbreak{} ，也都非零。四个矩阵均非奇异。\allowbreak{}\hskip.1em \RB{8.366}{7.994}{0.501}{%
\RG{0.000}{-0.001}{10.9091}{am}{3}%
}\hskip.1em\allowbreak{}

小结：将矩阵多项式的可逆性转化为相应标量多项式在特征值处是否为零。

题目来源：谢启鸿《复旦大学高等代数习题集（2023版）》，第93页。

% Original page 4.
\section*{2020S03}
设\allowbreak{}\hskip.1em \RB{6.729}{5.323}{0.619}{%
\RG{0.000}{0.001}{10.9091}{mi}{110}%
}\hskip.1em\allowbreak{} 阶复方阵

\[
\RB{123.694}{10.209}{69.316}{%
\RG{22.727}{9.283}{10.9092}{ex}{48}\RG{109.735}{9.283}{10.9092}{ex}{49}\RG{32.272}{0.002}{10.9091}{mi}{97}\RG{48.527}{0.002}{10.9091}{mi}{98}\RG{64.258}{0.002}{10.9091}{mi}{98}\RG{79.461}{0.002}{10.9091}{sy}{1}\RG{84.293}{0.002}{10.9091}{sy}{1}\RG{89.137}{0.002}{10.9091}{sy}{1}\RG{104.511}{0.002}{10.9091}{mi}{98}\RG{22.727}{-9.917}{10.9092}{ex}{66}\RG{109.735}{-9.917}{10.9092}{ex}{67}\RG{32.795}{-15.173}{10.9091}{mi}{99}\RG{47.992}{-15.173}{10.9091}{mi}{97}\RG{63.887}{-15.173}{10.9091}{rm}{48}\RG{79.461}{-15.173}{10.9091}{sy}{1}\RG{84.293}{-15.173}{10.9091}{sy}{1}\RG{89.137}{-15.173}{10.9091}{sy}{1}\RG{104.125}{-15.173}{10.9091}{rm}{48}\RG{22.727}{-16.462}{10.9092}{ex}{66}\RG{109.735}{-16.462}{10.9092}{ex}{67}\RG{22.727}{-23.008}{10.9092}{ex}{66}\RG{109.735}{-23.008}{10.9092}{ex}{67}\RG{22.727}{-29.553}{10.9092}{ex}{66}\RG{109.735}{-29.553}{10.9092}{ex}{67}\RG{32.795}{-30.348}{10.9091}{mi}{99}\RG{48.157}{-30.348}{10.9091}{rm}{48}\RG{63.732}{-30.348}{10.9091}{mi}{97}\RG{79.461}{-30.348}{10.9091}{sy}{1}\RG{84.293}{-30.348}{10.9091}{sy}{1}\RG{89.137}{-30.348}{10.9091}{sy}{1}\RG{104.125}{-30.348}{10.9091}{rm}{48}\RG{0.000}{-32.281}{10.9091}{mi}{65}\RG{11.212}{-32.281}{10.9091}{rm}{61}\RG{121.099}{-32.281}{10.9091}{mi}{58}\RG{22.727}{-36.099}{10.9092}{ex}{66}\RG{109.735}{-36.099}{10.9092}{ex}{67}\RT{33.639}{-42.034}{10.9091}{.}\RT{49.368}{-42.034}{10.9091}{.}\RT{65.098}{-42.034}{10.9091}{.}\RT{105.335}{-42.034}{10.9091}{.}\RG{22.727}{-42.644}{10.9092}{ex}{66}\RG{109.735}{-42.644}{10.9092}{ex}{67}\RT{80.972}{-43.030}{10.9091}{.}\RT{33.639}{-46.019}{10.9091}{.}\RT{49.368}{-46.019}{10.9091}{.}\RT{65.098}{-46.019}{10.9091}{.}\RT{85.217}{-46.019}{10.9091}{.}\RT{105.335}{-46.019}{10.9091}{.}\RT{89.462}{-49.008}{10.9091}{.}\RG{22.727}{-49.626}{10.9092}{ex}{64}\RG{109.735}{-49.626}{10.9092}{ex}{65}\RT{33.639}{-50.004}{10.9091}{.}\RT{49.368}{-50.004}{10.9091}{.}\RT{65.098}{-50.004}{10.9091}{.}\RT{105.335}{-50.004}{10.9091}{.}\RG{32.795}{-65.180}{10.9091}{mi}{99}\RG{48.157}{-65.180}{10.9091}{rm}{48}\RG{63.878}{-65.180}{10.9091}{rm}{48}\RG{79.461}{-65.180}{10.9091}{sy}{1}\RG{84.293}{-65.180}{10.9091}{sy}{1}\RG{89.137}{-65.180}{10.9091}{sy}{1}\RG{103.969}{-65.180}{10.9091}{mi}{97}%
}
\]
求\allowbreak{}\hskip.1em \RB{8.366}{8.311}{0.500}{%
\RG{0.000}{-0.000}{10.9091}{mi}{65}%
}\hskip.1em\allowbreak{} 可对角化的充要条件。

解：当\allowbreak{}\hskip.1em \RB{26.147}{7.764}{0.622}{%
\RG{0.000}{-0.002}{10.9091}{mi}{110}\RG{9.578}{-0.002}{10.9091}{rm}{61}\RG{21.076}{-0.002}{10.9091}{rm}{49}%
}\hskip.1em\allowbreak{} 时，\allowbreak{}\hskip.1em \RB{36.389}{8.680}{3.229}{%
\RG{0.000}{-0.002}{10.9091}{mi}{65}\RG{11.212}{-0.002}{10.9091}{rm}{61}\RG{22.710}{-0.002}{10.9091}{rm}{40}\RG{26.969}{-0.002}{10.9091}{mi}{97}\RG{32.736}{-0.002}{10.9091}{rm}{41}%
}\hskip.1em\allowbreak{} 总可对角化。以下设\allowbreak{}\hskip.1em \RB{26.491}{7.764}{1.997}{%
\RG{0.000}{-0.002}{10.9091}{mi}{110}\RG{9.578}{-0.002}{10.9091}{sy}{21}\RG{21.093}{-0.002}{10.9091}{rm}{50}%
}\hskip.1em\allowbreak{} ，结论为

\allowbreak{}\hskip.1em \RB{8.366}{8.306}{0.505}{%
\RG{0.000}{-0.005}{10.9091}{mi}{65}%
}\hskip.1em\allowbreak{} 可对角化\allowbreak{}\hskip.1em \RB{63.390}{8.306}{2.861}{%
\RG{0.000}{-0.005}{10.9091}{sy}{40}\RG{9.098}{-0.005}{10.9091}{sy}{41}\RG{33.940}{-0.005}{10.9091}{mi}{98}\RG{38.620}{-0.005}{10.9091}{mi}{99}\RG{46.373}{-0.005}{10.9091}{sy}{54}\RG{46.373}{-0.005}{10.9091}{rm}{61}\RG{57.872}{-0.005}{10.9091}{rm}{48}%
}\hskip.1em\allowbreak{} 或\allowbreak{}\hskip.1em \RB{46.543}{8.066}{0.745}{%
\RG{0.000}{-0.005}{10.9091}{mi}{98}\RG{7.712}{-0.005}{10.9091}{rm}{61}\RG{19.227}{-0.005}{10.9091}{mi}{99}\RG{26.979}{-0.005}{10.9091}{rm}{61}\RG{38.478}{-0.005}{10.9091}{rm}{48}\RG{43.948}{-0.005}{10.9091}{mi}{58}%
}\hskip.1em\allowbreak{}

令\allowbreak{}\hskip.1em \RB{61.955}{8.313}{2.214}{%
\RG{0.000}{0.002}{10.9091}{mi}{78}\RG{12.985}{0.002}{10.9091}{rm}{61}\RG{24.500}{0.002}{10.9091}{mi}{65}\RG{35.106}{0.002}{10.9091}{sy}{0}\RG{46.015}{0.002}{10.9091}{mi}{97}\RG{51.775}{0.002}{10.9091}{mi}{73}\RG{56.577}{-1.634}{7.9701}{mi}{110}%
}\hskip.1em\allowbreak{} ，\allowbreak{}\hskip.1em \RB{44.512}{5.324}{2.603}{%
\RG{0.000}{0.002}{10.9091}{mi}{101}\RG{5.080}{-1.634}{7.9701}{rm}{49}\RG{9.812}{0.002}{10.9091}{mi}{59}\RG{14.645}{0.002}{10.9091}{mi}{58}\RG{19.489}{0.002}{10.9091}{mi}{58}\RG{24.332}{0.002}{10.9091}{mi}{58}\RG{29.165}{0.002}{10.9091}{mi}{59}\RG{34.009}{0.002}{10.9091}{mi}{101}\RG{39.134}{-1.634}{7.9701}{mi}{110}%
}\hskip.1em\allowbreak{} 为标准基，并记

\[
\RB{287.366}{22.855}{17.396}{%
\RG{133.743}{22.365}{10.9092}{ex}{56}\RG{273.256}{22.366}{10.9092}{ex}{57}\RG{230.373}{13.638}{7.9701}{mi}{110}\RG{133.743}{12.547}{10.9092}{ex}{60}\RG{273.256}{12.547}{10.9092}{ex}{61}\RG{225.064}{10.365}{10.9092}{ex}{88}\RG{170.887}{4.506}{7.9701}{mi}{110}\RG{0.000}{0.002}{10.9091}{mi}{117}\RG{9.275}{0.002}{10.9091}{rm}{61}\RG{20.790}{0.002}{10.9091}{mi}{101}\RG{33.027}{0.002}{10.9091}{rm}{43}\RG{43.936}{0.002}{10.9091}{sy}{1}\RG{48.769}{0.002}{10.9091}{sy}{1}\RG{53.613}{0.002}{10.9091}{sy}{1}\RG{59.087}{0.002}{10.9091}{rm}{43}\RG{69.997}{0.002}{10.9091}{mi}{101}\RG{80.713}{0.002}{10.9091}{mi}{59}\RG{107.364}{0.002}{10.9091}{mi}{87}\RG{122.228}{0.002}{10.9091}{rm}{61}\RG{143.440}{0.002}{10.9091}{mi}{120}\RG{152.705}{0.002}{10.9091}{sy}{50}\RG{163.008}{0.002}{10.9091}{bb}{67}\RG{179.554}{0.002}{10.9091}{rm}{58}\RG{185.615}{0.002}{10.9091}{mi}{120}\RG{199.613}{0.002}{10.9091}{rm}{61}\RG{211.111}{0.002}{10.9091}{rm}{48}\RG{216.582}{0.002}{10.9091}{mi}{59}\RG{242.639}{0.002}{10.9091}{mi}{120}\RG{256.287}{0.002}{10.9091}{rm}{61}\RG{267.785}{0.002}{10.9091}{rm}{48}\RG{284.771}{0.002}{10.9091}{mi}{58}\RG{25.870}{-1.635}{7.9701}{rm}{50}\RG{75.077}{-1.635}{7.9701}{mi}{110}\RG{191.850}{-1.635}{7.9701}{rm}{49}\RG{248.874}{-1.635}{7.9701}{mi}{106}\RG{133.743}{-7.089}{10.9092}{ex}{58}\RG{273.256}{-7.089}{10.9092}{ex}{59}\RG{225.590}{-13.085}{7.9701}{mi}{106}\RG{229.474}{-13.085}{7.9701}{rm}{61}\RG{236.057}{-13.085}{7.9701}{rm}{50}%
}
\]
两个线性条件相互独立，故dim\allowbreak{}\hskip.1em \RB{51.783}{7.950}{0.741}{%
\RG{0.000}{-0.001}{10.9091}{mi}{87}\RG{14.898}{-0.001}{10.9091}{rm}{61}\RG{26.463}{-0.001}{10.9091}{mi}{110}\RG{35.455}{-0.001}{10.9091}{sy}{0}\RG{46.384}{-0.001}{10.9091}{rm}{50}%
}\hskip.1em\allowbreak{} 。若\allowbreak{}\hskip.1em \RB{68.560}{8.201}{2.617}{%
\RG{0.000}{-0.001}{10.9091}{mi}{11}\RG{7.004}{-0.001}{10.9091}{mi}{101}\RG{12.098}{-1.637}{7.9701}{rm}{49}\RG{19.275}{-0.001}{10.9091}{rm}{43}\RG{30.204}{-0.001}{10.9091}{mi}{12}\RG{36.935}{-0.001}{10.9091}{mi}{117}\RG{46.274}{-0.001}{10.9091}{sy}{50}\RG{56.627}{-0.001}{10.9091}{mi}{87}%
}\hskip.1em\allowbreak{} ，则\allowbreak{}\hskip.1em \RB{27.169}{7.765}{0.741}{%
\RG{0.000}{-0.001}{10.9091}{mi}{11}\RG{10.098}{-0.001}{10.9091}{rm}{61}\RG{21.651}{-0.001}{10.9091}{rm}{48}%
}\hskip.1em\allowbreak{} 且\allowbreak{}\hskip.1em \RB{60.758}{8.681}{3.228}{%
\RG{0.000}{-0.001}{10.9091}{rm}{40}\RG{4.243}{-0.001}{10.9091}{mi}{110}\RG{13.235}{-0.001}{10.9091}{sy}{0}\RG{24.164}{-0.001}{10.9091}{rm}{49}\RG{29.619}{-0.001}{10.9091}{rm}{41}\RG{33.861}{-0.001}{10.9091}{mi}{12}\RG{43.687}{-0.001}{10.9091}{rm}{61}\RG{55.240}{-0.001}{10.9091}{rm}{48}%
}\hskip.1em\allowbreak{} ，从而\allowbreak{}\hskip.1em \RB{26.793}{8.206}{2.612}{%
\RG{0.000}{0.004}{10.9091}{mi}{12}\RG{9.776}{0.004}{10.9091}{rm}{61}\RG{21.274}{0.004}{10.9091}{rm}{48}%
}\hskip.1em\allowbreak{} 。因此

\[
\RB{109.505}{9.292}{3.228}{%
\RG{7.878}{4.503}{7.9701}{mi}{110}\RG{0.000}{-0.001}{10.9091}{bb}{67}\RG{16.545}{-0.001}{10.9091}{rm}{61}\RG{28.060}{-0.001}{10.9091}{mi}{87}\RG{42.303}{-0.001}{10.9091}{sy}{8}\RT{53.212}{-0.001}{10.9091}{s}\RT{57.511}{-0.001}{10.9091}{p}\RT{63.576}{-0.001}{10.9091}{a}\RT{69.031}{-0.001}{10.9091}{n}\RG{75.096}{-0.001}{10.9091}{sy}{102}\RG{80.551}{-0.001}{10.9091}{mi}{101}\RG{90.363}{-0.001}{10.9091}{mi}{59}\RG{95.196}{-0.001}{10.9091}{mi}{117}\RG{101.456}{-0.001}{10.9091}{sy}{103}\RG{106.911}{-0.001}{10.9091}{mi}{58}\RG{85.630}{-1.637}{7.9701}{rm}{49}%
}
\]
直接计算得

\[
\RB{222.915}{8.685}{3.224}{%
\RG{0.000}{0.003}{10.9091}{mi}{78}\RG{9.955}{0.003}{10.9091}{sy}{106}\RG{25.679}{0.003}{10.9091}{rm}{61}\RG{37.178}{0.003}{10.9091}{rm}{48}\RG{42.649}{0.003}{10.9091}{mi}{59}\RG{69.301}{0.003}{10.9091}{mi}{78}\RG{79.250}{0.003}{10.9091}{mi}{101}\RG{92.113}{0.003}{10.9091}{rm}{61}\RG{103.628}{0.003}{10.9091}{mi}{99}\RG{108.341}{0.003}{10.9091}{mi}{117}\RG{114.581}{0.003}{10.9091}{mi}{59}\RG{141.232}{0.003}{10.9091}{mi}{78}\RG{151.181}{0.003}{10.9091}{mi}{117}\RG{160.491}{0.003}{10.9091}{rm}{61}\RG{171.990}{0.003}{10.9091}{rm}{40}\RG{176.248}{0.003}{10.9091}{mi}{110}\RG{185.221}{0.003}{10.9091}{sy}{0}\RG{196.130}{0.003}{10.9091}{rm}{49}\RG{201.584}{0.003}{10.9091}{rm}{41}\RG{205.827}{0.003}{10.9091}{mi}{98}\RG{210.507}{0.003}{10.9091}{mi}{101}\RG{220.321}{0.003}{10.9091}{mi}{58}\RG{84.351}{-1.633}{7.9701}{rm}{49}\RG{215.589}{-1.633}{7.9701}{rm}{49}\RG{12.985}{-1.685}{7.9701}{mi}{87}%
}
\]
所以在与此分解相适应的基下，

\[
\RB{139.118}{12.313}{21.404}{%
\RG{63.079}{11.376}{10.9092}{ex}{32}\RG{125.948}{11.376}{10.9092}{ex}{33}\RG{71.716}{0.001}{10.9091}{mi}{97}\RG{87.446}{0.001}{10.9091}{rm}{40}\RG{91.688}{0.001}{10.9091}{mi}{110}\RG{100.660}{0.001}{10.9091}{sy}{0}\RG{111.569}{0.001}{10.9091}{rm}{49}\RG{117.024}{0.001}{10.9091}{rm}{41}\RG{121.266}{0.001}{10.9091}{mi}{98}\RG{0.000}{-7.279}{10.9091}{mi}{65}\RG{11.212}{-7.279}{10.9091}{sy}{24}\RG{22.727}{-7.279}{10.9091}{mi}{97}\RG{28.487}{-7.279}{10.9091}{mi}{73}\RG{52.170}{-7.279}{10.9091}{sy}{8}\RG{136.403}{-7.279}{10.9091}{mi}{59}\RG{33.289}{-8.915}{7.9701}{mi}{110}\RG{38.427}{-8.915}{7.9701}{sy}{0}\RG{45.013}{-8.915}{7.9701}{rm}{50}\RG{72.239}{-15.174}{10.9091}{mi}{99}\RG{103.799}{-15.174}{10.9091}{mi}{97}%
}
\]
其中\allowbreak{}\hskip.1em \RB{26.475}{7.765}{0.621}{%
\RG{0.000}{-0.001}{10.9091}{mi}{110}\RG{9.578}{-0.001}{10.9091}{rm}{61}\RG{21.076}{-0.001}{10.9091}{rm}{50}%
}\hskip.1em\allowbreak{} 时省略第一个直和项。

若\allowbreak{}\hskip.1em \RB{30.912}{8.306}{2.861}{%
\RG{0.000}{-0.005}{10.9091}{mi}{98}\RG{4.680}{-0.005}{10.9091}{mi}{99}\RG{13.164}{-0.005}{10.9091}{sy}{54}\RG{13.164}{-0.005}{10.9091}{rm}{61}\RG{25.393}{-0.005}{10.9091}{rm}{48}%
}\hskip.1em\allowbreak{} ，取\allowbreak{}\hskip.1em \RB{28.029}{8.252}{0.941}{%
\RG{0.000}{-0.005}{10.9091}{mi}{14}\RG{9.022}{-0.005}{10.9091}{sy}{50}\RG{20.056}{-0.005}{10.9091}{bb}{67}%
}\hskip.1em\allowbreak{} 满足\allowbreak{}\hskip.1em \RB{70.276}{9.754}{3.232}{%
\RG{0.000}{-0.005}{10.9091}{mi}{14}\RG{5.261}{3.954}{7.9701}{rm}{50}\RG{13.754}{-0.005}{10.9091}{rm}{61}\RG{25.983}{-0.005}{10.9091}{rm}{40}\RG{30.243}{-0.005}{10.9091}{mi}{110}\RG{39.507}{-0.005}{10.9091}{sy}{0}\RG{50.708}{-0.005}{10.9091}{rm}{49}\RG{56.163}{-0.005}{10.9091}{rm}{41}\RG{60.406}{-0.005}{10.9091}{mi}{98}\RG{65.086}{-0.005}{10.9091}{mi}{99}%
}\hskip.1em\allowbreak{} 。此时\allowbreak{}\hskip.1em \RB{26.770}{8.306}{2.861}{%
\RG{0.000}{-0.005}{10.9091}{mi}{14}\RG{9.022}{-0.005}{10.9091}{sy}{54}\RG{9.022}{-0.005}{10.9091}{rm}{61}\RG{21.251}{-0.005}{10.9091}{rm}{48}%
}\hskip.1em\allowbreak{} ，右侧二阶矩阵有两个不同的特征值\allowbreak{}\hskip.1em \RB{24.531}{8.256}{1.405}{%
\RG{0.000}{0.000}{10.9091}{mi}{97}\RG{8.191}{0.000}{10.9091}{rm}{43}\RG{19.100}{0.000}{10.9091}{mi}{14}%
}\hskip.1em\allowbreak{} 与\allowbreak{}\hskip.1em \RB{24.530}{8.256}{0.631}{%
\RG{0.000}{0.000}{10.9091}{mi}{97}\RG{8.190}{0.000}{10.9091}{sy}{0}\RG{19.099}{0.000}{10.9091}{mi}{14}%
}\hskip.1em\allowbreak{} ，故可对角化。于是\allowbreak{}\hskip.1em \RB{8.366}{8.311}{0.500}{%
\RG{0.000}{0.000}{10.9091}{mi}{65}%
}\hskip.1em\allowbreak{} 可对角化；其特征值为\allowbreak{}\hskip.1em \RB{5.933}{5.322}{0.620}{%
\RG{0.000}{0.000}{10.9091}{mi}{97}%
}\hskip.1em\allowbreak{} （\allowbreak{}\hskip.1em \RB{25.279}{7.766}{0.620}{%
\RG{0.000}{0.000}{10.9091}{mi}{110}\RG{8.972}{0.000}{10.9091}{sy}{0}\RG{19.881}{0.000}{10.9091}{rm}{50}%
}\hskip.1em\allowbreak{} 重）以及\allowbreak{}\hskip.1em \RB{24.531}{8.256}{0.631}{%
\RG{0.000}{0.000}{10.9091}{mi}{97}\RG{8.191}{0.000}{10.9091}{sy}{6}\RG{19.100}{0.000}{10.9091}{mi}{14}%
}\hskip.1em\allowbreak{} 。

若\allowbreak{}\hskip.1em \RB{43.995}{8.068}{0.744}{%
\RG{0.000}{-0.003}{10.9091}{mi}{98}\RG{7.712}{-0.004}{10.9091}{rm}{61}\RG{19.227}{-0.003}{10.9091}{mi}{99}\RG{26.978}{-0.004}{10.9091}{rm}{61}\RG{38.477}{-0.004}{10.9091}{rm}{48}%
}\hskip.1em\allowbreak{} ，则\allowbreak{}\hskip.1em \RB{38.667}{8.307}{2.220}{%
\RG{0.000}{-0.004}{10.9091}{mi}{65}\RG{11.212}{-0.003}{10.9091}{rm}{61}\RG{22.727}{-0.004}{10.9091}{mi}{97}\RG{28.487}{-0.004}{10.9091}{mi}{73}\RG{33.289}{-1.640}{7.9701}{mi}{110}%
}\hskip.1em\allowbreak{} ，也可对角化。若\allowbreak{}\hskip.1em \RB{29.578}{8.074}{0.737}{%
\RG{0.000}{0.003}{10.9091}{mi}{98}\RG{4.680}{0.003}{10.9091}{mi}{99}\RG{12.496}{0.003}{10.9091}{rm}{61}\RG{24.060}{0.003}{10.9091}{rm}{48}%
}\hskip.1em\allowbreak{} 而\allowbreak{}\hskip.1em \RB{61.061}{8.685}{3.224}{%
\RG{0.000}{0.003}{10.9091}{rm}{40}\RG{4.242}{0.003}{10.9091}{mi}{98}\RG{8.922}{0.003}{10.9091}{mi}{59}\RG{13.755}{0.003}{10.9091}{mi}{99}\RG{18.493}{0.003}{10.9091}{rm}{41}\RG{25.829}{0.003}{10.9091}{sy}{54}\RG{25.829}{0.003}{10.9091}{rm}{61}\RG{37.393}{0.003}{10.9091}{rm}{40}\RG{41.626}{0.003}{10.9091}{rm}{48}\RG{47.105}{0.003}{10.9091}{mi}{59}\RG{51.953}{0.003}{10.9091}{rm}{48}\RG{57.408}{0.003}{10.9091}{rm}{41}%
}\hskip.1em\allowbreak{} ，则\allowbreak{}\hskip.1em \RB{30.130}{8.314}{2.853}{%
\RG{0.000}{0.003}{10.9091}{mi}{78}\RG{13.048}{0.003}{10.9091}{sy}{54}\RG{13.048}{0.003}{10.9091}{rm}{61}\RG{24.612}{0.003}{10.9091}{rm}{48}%
}\hskip.1em\allowbreak{} ，且由上述直和分解可得\allowbreak{}\hskip.1em \RB{34.862}{9.761}{0.737}{%
\RG{0.000}{0.003}{10.9091}{mi}{78}\RG{9.954}{3.961}{7.9701}{rm}{50}\RG{17.780}{0.003}{10.9091}{rm}{61}\RG{29.344}{0.003}{10.9091}{rm}{48}%
}\hskip.1em\allowbreak{} 。因此\allowbreak{}\hskip.1em \RB{10.111}{7.954}{0.500}{%
\RG{0.000}{0.003}{10.9091}{mi}{78}%
}\hskip.1em\allowbreak{} 的特征值全为\allowbreak{}\hskip.1em \RB{5.518}{7.764}{0.742}{%
\RG{0.000}{-0.002}{10.9091}{rm}{48}%
}\hskip.1em\allowbreak{} 。若\allowbreak{}\hskip.1em \RB{8.366}{8.309}{0.502}{%
\RG{0.000}{-0.002}{10.9091}{mi}{65}%
}\hskip.1em\allowbreak{} 可对角化，则\allowbreak{}\hskip.1em \RB{61.414}{8.309}{2.219}{%
\RG{0.000}{-0.002}{10.9091}{mi}{78}\RG{12.985}{-0.002}{10.9091}{rm}{61}\RG{24.500}{-0.002}{10.9091}{mi}{65}\RG{34.835}{-0.002}{10.9091}{sy}{0}\RG{45.474}{-0.002}{10.9091}{mi}{97}\RG{51.234}{-0.002}{10.9091}{mi}{73}\RG{56.036}{-1.639}{7.9701}{mi}{110}%
}\hskip.1em\allowbreak{} 也可对角化，只能等于零矩阵，与\allowbreak{}\hskip.1em \RB{30.002}{8.309}{2.858}{%
\RG{0.000}{-0.002}{10.9091}{mi}{78}\RG{12.985}{-0.002}{10.9091}{sy}{54}\RG{12.985}{-0.002}{10.9091}{rm}{61}\RG{24.483}{-0.002}{10.9091}{rm}{48}%
}\hskip.1em\allowbreak{} 矛盾。充要条件得证。\allowbreak{}\hskip.1em \RB{8.366}{7.998}{0.500}{%
\RG{0.000}{0.003}{10.9091}{am}{3}%
}\hskip.1em\allowbreak{}

小结：把空间分成一个零作用子空间与一个二维不变子空间，便可完整处理参数退化的情形。

题目来源：谢启鸿《复旦大学高等代数习题集（2023版）》，第93页。

% Original page 5.
\section*{2020S04}
设\allowbreak{}\hskip.1em \RB{55.675}{8.683}{3.226}{%
\RG{0.000}{0.001}{10.9091}{mi}{65}\RG{11.212}{0.001}{10.9091}{sy}{50}\RG{21.515}{0.001}{10.9091}{mi}{77}\RG{32.098}{-1.635}{7.9701}{mi}{110}\RG{37.735}{0.001}{10.9091}{rm}{40}\RG{41.977}{0.001}{10.9091}{mi}{75}\RG{52.022}{0.001}{10.9091}{rm}{41}%
}\hskip.1em\allowbreak{} ，\allowbreak{}\hskip.1em \RB{81.277}{8.683}{3.226}{%
\RG{0.000}{0.001}{10.9091}{mi}{102}\RG{6.515}{0.001}{10.9091}{rm}{40}\RG{10.758}{0.001}{10.9091}{mi}{120}\RG{16.992}{0.001}{10.9091}{rm}{41}\RG{21.235}{0.001}{10.9091}{mi}{59}\RG{26.068}{0.001}{10.9091}{mi}{103}\RG{31.678}{0.001}{10.9091}{rm}{40}\RG{35.920}{0.001}{10.9091}{mi}{120}\RG{42.155}{0.001}{10.9091}{rm}{41}\RG{49.428}{0.001}{10.9091}{sy}{50}\RG{59.731}{0.001}{10.9091}{mi}{75}\RG{69.777}{0.001}{10.9091}{rm}{91}\RG{72.808}{0.001}{10.9091}{mi}{120}\RG{79.043}{0.001}{10.9091}{rm}{93}%
}\hskip.1em\allowbreak{} 互素，且它们在\allowbreak{}\hskip.1em \RB{7.973}{8.181}{0.706}{%
\RG{0.000}{0.001}{10.9091}{bb}{67}%
}\hskip.1em\allowbreak{} 中均无重根。证明：若

\[
\RB{129.332}{9.290}{3.230}{%
\RT{0.000}{-0.003}{10.9091}{r}\RT{4.276}{-0.003}{10.9091}{a}\RT{9.731}{-0.003}{10.9091}{n}\RT{15.797}{-0.003}{10.9091}{k}\RG{23.375}{-0.003}{10.9091}{mi}{102}\RG{29.890}{-0.003}{10.9091}{rm}{40}\RG{34.132}{-0.003}{10.9091}{mi}{65}\RG{42.314}{-0.003}{10.9091}{rm}{41}\RG{48.969}{-0.003}{10.9091}{rm}{43}\RT{59.891}{-0.003}{10.9091}{r}\RT{64.167}{-0.003}{10.9091}{a}\RT{69.621}{-0.003}{10.9091}{n}\RT{75.687}{-0.003}{10.9091}{k}\RG{83.265}{-0.003}{10.9091}{mi}{103}\RG{88.860}{-0.003}{10.9091}{rm}{40}\RG{93.103}{-0.003}{10.9091}{mi}{65}\RG{101.284}{-0.003}{10.9091}{rm}{41}\RG{108.538}{-0.003}{10.9091}{rm}{61}\RG{120.072}{-0.003}{10.9091}{mi}{110}\RG{126.618}{-0.003}{10.9091}{mi}{59}%
}
\]
则\allowbreak{}\hskip.1em \RB{8.366}{8.314}{0.500}{%
\RG{0.000}{0.003}{10.9091}{mi}{65}%
}\hskip.1em\allowbreak{} 在复数域上可对角化。

证明：在\allowbreak{}\hskip.1em \RB{13.257}{8.182}{0.705}{%
\RG{0.000}{0.002}{10.9091}{bb}{67}\RG{7.879}{3.961}{7.9701}{mi}{110}%
}\hskip.1em\allowbreak{} 上考虑\allowbreak{}\hskip.1em \RB{8.366}{8.313}{0.500}{%
\RG{0.000}{0.002}{10.9091}{mi}{65}%
}\hskip.1em\allowbreak{} ；扩充数域不改变矩阵的秩，因为秩由非零子式的最大阶数确定。由秩与零度之和等于\allowbreak{}\hskip.1em \RB{6.729}{5.319}{0.623}{%
\RG{0.000}{-0.003}{10.9091}{mi}{110}%
}\hskip.1em\allowbreak{} ，有

\[
\RB{155.866}{9.295}{3.225}{%
\RT{0.000}{0.002}{10.9091}{d}\RT{6.066}{0.002}{10.9091}{i}\RT{9.098}{0.002}{10.9091}{m}\RT{19.997}{0.002}{10.9091}{k}\RT{25.757}{0.002}{10.9091}{e}\RT{30.600}{0.002}{10.9091}{r}\RG{36.702}{0.002}{10.9091}{mi}{102}\RG{43.217}{0.002}{10.9091}{rm}{40}\RG{47.459}{0.002}{10.9091}{mi}{65}\RG{55.641}{0.002}{10.9091}{rm}{41}\RG{62.296}{0.002}{10.9091}{rm}{43}\RT{73.218}{0.002}{10.9091}{d}\RT{79.283}{0.002}{10.9091}{i}\RT{82.316}{0.002}{10.9091}{m}\RT{93.214}{0.002}{10.9091}{k}\RT{98.974}{0.002}{10.9091}{e}\RT{103.818}{0.002}{10.9091}{r}\RG{109.920}{0.002}{10.9091}{mi}{103}\RG{115.514}{0.002}{10.9091}{rm}{40}\RG{119.757}{0.002}{10.9091}{mi}{65}\RG{127.939}{0.002}{10.9091}{rm}{41}\RG{135.194}{0.002}{10.9091}{rm}{61}\RG{146.726}{0.002}{10.9091}{mi}{110}\RG{153.272}{0.002}{10.9091}{mi}{58}%
}
\]
由\allowbreak{}\hskip.1em \RB{16.406}{8.192}{2.735}{%
\RG{0.000}{0.001}{10.9091}{mi}{102}\RG{5.902}{0.001}{10.9091}{mi}{59}\RG{10.735}{0.001}{10.9091}{mi}{103}%
}\hskip.1em\allowbreak{} 互素，存在\allowbreak{}\hskip.1em \RB{51.652}{8.683}{3.226}{%
\RG{0.000}{0.001}{10.9091}{mi}{117}\RG{6.240}{0.001}{10.9091}{mi}{59}\RG{11.073}{0.001}{10.9091}{mi}{118}\RG{19.803}{0.001}{10.9091}{sy}{50}\RG{30.106}{0.001}{10.9091}{mi}{75}\RG{40.151}{0.001}{10.9091}{rm}{91}\RG{43.182}{0.001}{10.9091}{mi}{120}\RG{49.417}{0.001}{10.9091}{rm}{93}%
}\hskip.1em\allowbreak{} ，使\allowbreak{}\hskip.1em \RB{56.968}{8.192}{2.735}{%
\RG{0.000}{0.001}{10.9091}{mi}{117}\RG{6.240}{0.001}{10.9091}{mi}{102}\RG{15.184}{0.001}{10.9091}{rm}{43}\RG{26.093}{0.001}{10.9091}{mi}{118}\RG{31.755}{0.001}{10.9091}{mi}{103}\RG{40.398}{0.001}{10.9091}{rm}{61}\RG{51.897}{0.001}{10.9091}{rm}{49}%
}\hskip.1em\allowbreak{} 。若\allowbreak{}\hskip.1em \RB{15.443}{6.392}{0.935}{%
\RG{0.000}{0.001}{10.9091}{mi}{122}\RG{8.583}{0.001}{10.9091}{sy}{50}%
}\hskip.1em\allowbreak{} ker\allowbreak{}\hskip.1em \RB{32.772}{8.683}{3.226}{%
\RG{0.000}{0.001}{10.9091}{mi}{102}\RG{6.515}{0.001}{10.9091}{rm}{40}\RG{10.757}{0.001}{10.9091}{mi}{65}\RG{18.939}{0.001}{10.9091}{rm}{41}\RG{25.606}{0.001}{10.9091}{sy}{92}%
}\hskip.1em\allowbreak{} ker\allowbreak{}\hskip.1em \RB{21.672}{8.683}{3.226}{%
\RG{0.000}{0.001}{10.9091}{mi}{103}\RG{5.594}{0.001}{10.9091}{rm}{40}\RG{9.837}{0.001}{10.9091}{mi}{65}\RG{18.019}{0.001}{10.9091}{rm}{41}%
}\hskip.1em\allowbreak{} ，则

\[
\RB{157.835}{8.679}{3.230}{%
\RG{0.000}{-0.003}{10.9091}{mi}{122}\RG{8.584}{-0.003}{10.9091}{rm}{61}\RG{20.099}{-0.003}{10.9091}{mi}{117}\RG{26.344}{-0.003}{10.9091}{rm}{40}\RG{30.586}{-0.003}{10.9091}{mi}{65}\RG{38.768}{-0.003}{10.9091}{rm}{41}\RG{43.011}{-0.003}{10.9091}{mi}{102}\RG{49.526}{-0.003}{10.9091}{rm}{40}\RG{53.768}{-0.003}{10.9091}{mi}{65}\RG{61.951}{-0.003}{10.9091}{rm}{41}\RG{66.193}{-0.003}{10.9091}{mi}{122}\RG{74.171}{-0.003}{10.9091}{rm}{43}\RG{85.080}{-0.003}{10.9091}{mi}{118}\RG{90.759}{-0.003}{10.9091}{rm}{40}\RG{95.002}{-0.003}{10.9091}{mi}{65}\RG{103.183}{-0.003}{10.9091}{rm}{41}\RG{107.426}{-0.003}{10.9091}{mi}{103}\RG{113.021}{-0.003}{10.9091}{rm}{40}\RG{117.263}{-0.003}{10.9091}{mi}{65}\RG{125.445}{-0.003}{10.9091}{rm}{41}\RG{129.687}{-0.003}{10.9091}{mi}{122}\RG{138.271}{-0.003}{10.9091}{rm}{61}\RG{149.769}{-0.003}{10.9091}{rm}{48}\RG{155.240}{-0.003}{10.9091}{mi}{58}%
}
\]
于是

\[
\RB{122.829}{9.291}{3.229}{%
\RG{7.879}{4.503}{7.9701}{mi}{110}\RG{0.000}{-0.002}{10.9091}{bb}{67}\RG{16.545}{-0.002}{10.9091}{rm}{61}\RT{28.061}{-0.002}{10.9091}{k}\RT{33.821}{-0.002}{10.9091}{e}\RT{38.665}{-0.002}{10.9091}{r}\RG{44.759}{-0.002}{10.9091}{mi}{102}\RG{51.274}{-0.002}{10.9091}{rm}{40}\RG{55.516}{-0.002}{10.9091}{mi}{65}\RG{63.699}{-0.002}{10.9091}{rm}{41}\RG{70.366}{-0.002}{10.9091}{sy}{8}\RT{81.275}{-0.002}{10.9091}{k}\RT{87.035}{-0.002}{10.9091}{e}\RT{91.878}{-0.002}{10.9091}{r}\RG{97.973}{-0.002}{10.9091}{mi}{103}\RG{103.568}{-0.002}{10.9091}{rm}{40}\RG{107.810}{-0.002}{10.9091}{mi}{65}\RG{115.992}{-0.002}{10.9091}{rm}{41}\RG{120.234}{-0.002}{10.9091}{mi}{58}%
}
\]
令\allowbreak{}\hskip.1em \RB{33.004}{8.190}{2.737}{%
\RG{0.000}{-0.001}{10.9091}{mi}{104}\RG{9.316}{-0.001}{10.9091}{rm}{61}\RG{20.831}{-0.001}{10.9091}{mi}{102}\RG{27.333}{-0.001}{10.9091}{mi}{103}%
}\hskip.1em\allowbreak{} 。对任意\allowbreak{}\hskip.1em \RB{52.262}{6.859}{2.138}{%
\RG{0.000}{-0.001}{10.9091}{mi}{122}\RG{8.584}{-0.001}{10.9091}{rm}{61}\RG{20.099}{-0.001}{10.9091}{mi}{122}\RG{25.172}{-1.638}{7.9701}{rm}{49}\RG{32.153}{-0.001}{10.9091}{rm}{43}\RG{42.886}{-0.001}{10.9091}{mi}{122}\RG{47.960}{-1.638}{7.9701}{rm}{50}%
}\hskip.1em\allowbreak{} ，其中\allowbreak{}\hskip.1em \RB{19.696}{6.390}{2.138}{%
\RG{0.000}{-0.001}{10.9091}{mi}{122}\RG{5.073}{-1.638}{7.9701}{rm}{49}\RG{12.836}{-0.001}{10.9091}{sy}{50}%
}\hskip.1em\allowbreak{} ker\allowbreak{}\hskip.1em \RB{22.592}{8.681}{3.228}{%
\RG{0.000}{-0.001}{10.9091}{mi}{102}\RG{6.515}{-0.001}{10.9091}{rm}{40}\RG{10.758}{-0.001}{10.9091}{mi}{65}\RG{18.939}{-0.001}{10.9091}{rm}{41}%
}\hskip.1em\allowbreak{} 、\allowbreak{}\hskip.1em \RB{19.696}{6.390}{2.138}{%
\RG{0.000}{-0.001}{10.9091}{mi}{122}\RG{5.073}{-1.638}{7.9701}{rm}{50}\RG{12.836}{-0.001}{10.9091}{sy}{50}%
}\hskip.1em\allowbreak{} ker\allowbreak{}\hskip.1em \RB{21.672}{8.681}{3.228}{%
\RG{0.000}{-0.001}{10.9091}{mi}{103}\RG{5.595}{-0.001}{10.9091}{rm}{40}\RG{9.837}{-0.001}{10.9091}{mi}{65}\RG{18.019}{-0.001}{10.9091}{rm}{41}%
}\hskip.1em\allowbreak{} ，利用\allowbreak{}\hskip.1em \RB{22.593}{8.681}{3.228}{%
\RG{0.000}{-0.001}{10.9091}{mi}{102}\RG{6.516}{-0.001}{10.9091}{rm}{40}\RG{10.758}{-0.001}{10.9091}{mi}{65}\RG{18.940}{-0.001}{10.9091}{rm}{41}%
}\hskip.1em\allowbreak{} 与\allowbreak{}\hskip.1em \RB{21.672}{8.681}{3.228}{%
\RG{0.000}{-0.001}{10.9091}{mi}{103}\RG{5.595}{-0.001}{10.9091}{rm}{40}\RG{9.838}{-0.001}{10.9091}{mi}{65}\RG{18.019}{-0.001}{10.9091}{rm}{41}%
}\hskip.1em\allowbreak{} 可交换，得到

\[
\RB{189.478}{8.681}{3.228}{%
\RG{0.000}{-0.001}{10.9091}{mi}{104}\RG{6.285}{-0.001}{10.9091}{rm}{40}\RG{10.528}{-0.001}{10.9091}{mi}{65}\RG{18.710}{-0.001}{10.9091}{rm}{41}\RG{22.952}{-0.001}{10.9091}{mi}{122}\RG{31.536}{-0.001}{10.9091}{rm}{61}\RG{43.051}{-0.001}{10.9091}{mi}{103}\RG{48.645}{-0.001}{10.9091}{rm}{40}\RG{52.888}{-0.001}{10.9091}{mi}{65}\RG{61.070}{-0.001}{10.9091}{rm}{41}\RG{65.313}{-0.001}{10.9091}{mi}{102}\RG{71.828}{-0.001}{10.9091}{rm}{40}\RG{76.071}{-0.001}{10.9091}{mi}{65}\RG{84.253}{-0.001}{10.9091}{rm}{41}\RG{88.495}{-0.001}{10.9091}{mi}{122}\RG{100.725}{-0.001}{10.9091}{rm}{43}\RG{111.634}{-0.001}{10.9091}{mi}{102}\RG{118.149}{-0.001}{10.9091}{rm}{40}\RG{122.391}{-0.001}{10.9091}{mi}{65}\RG{130.573}{-0.001}{10.9091}{rm}{41}\RG{134.816}{-0.001}{10.9091}{mi}{103}\RG{140.410}{-0.001}{10.9091}{rm}{40}\RG{144.653}{-0.001}{10.9091}{mi}{65}\RG{152.835}{-0.001}{10.9091}{rm}{41}\RG{157.077}{-0.001}{10.9091}{mi}{122}\RG{169.913}{-0.001}{10.9091}{rm}{61}\RG{181.412}{-0.001}{10.9091}{rm}{48}\RG{186.884}{-0.001}{10.9091}{mi}{58}\RG{93.568}{-1.638}{7.9701}{rm}{49}\RG{162.150}{-1.638}{7.9701}{rm}{50}%
}
\]
故\allowbreak{}\hskip.1em \RB{42.992}{8.680}{3.229}{%
\RG{0.000}{-0.002}{10.9091}{mi}{104}\RG{6.285}{-0.002}{10.9091}{rm}{40}\RG{10.528}{-0.002}{10.9091}{mi}{65}\RG{18.710}{-0.002}{10.9091}{rm}{41}\RG{25.965}{-0.002}{10.9091}{rm}{61}\RG{37.474}{-0.002}{10.9091}{rm}{48}%
}\hskip.1em\allowbreak{} 。又因\allowbreak{}\hskip.1em \RB{16.406}{8.189}{2.738}{%
\RG{0.000}{-0.002}{10.9091}{mi}{102}\RG{5.902}{-0.002}{10.9091}{mi}{59}\RG{10.735}{-0.002}{10.9091}{mi}{103}%
}\hskip.1em\allowbreak{} 各自无重根且无公共根，\allowbreak{}\hskip.1em \RB{6.456}{8.069}{0.622}{%
\RG{0.000}{-0.002}{10.9091}{mi}{104}%
}\hskip.1em\allowbreak{} 也无重根。

下面直接说明存在无重根的零化多项式为何推出可对角化。将\allowbreak{}\hskip.1em \RB{6.456}{8.075}{0.616}{%
\RG{0.000}{0.004}{10.9091}{mi}{104}%
}\hskip.1em\allowbreak{} 写为

\[
\RB{149.696}{17.654}{15.210}{%
\RG{47.413}{13.639}{7.9701}{mi}{114}\RG{42.472}{10.366}{10.9092}{ex}{89}\RG{0.000}{0.003}{10.9091}{mi}{104}\RG{6.285}{0.003}{10.9091}{rm}{40}\RG{10.527}{0.003}{10.9091}{mi}{120}\RG{16.762}{0.003}{10.9091}{rm}{41}\RG{24.017}{0.003}{10.9091}{rm}{61}\RG{35.550}{0.003}{10.9091}{mi}{99}\RG{56.794}{0.003}{10.9091}{rm}{40}\RG{61.037}{0.003}{10.9091}{mi}{120}\RG{69.696}{0.003}{10.9091}{sy}{0}\RG{80.605}{0.003}{10.9091}{mi}{21}\RG{91.351}{0.003}{10.9091}{rm}{41}\RG{95.594}{0.003}{10.9091}{mi}{59}\RG{122.245}{0.003}{10.9091}{mi}{99}\RG{130.012}{0.003}{10.9091}{sy}{54}\RG{130.012}{0.003}{10.9091}{rm}{61}\RG{141.510}{0.003}{10.9091}{rm}{48}\RG{146.981}{0.003}{10.9091}{mi}{59}\RG{86.969}{-1.634}{7.9701}{mi}{106}\RG{42.089}{-13.084}{7.9701}{mi}{106}\RG{45.973}{-13.084}{7.9701}{rm}{61}\RG{52.557}{-13.084}{7.9701}{rm}{49}%
}
\]
其中\allowbreak{}\hskip.1em \RB{45.885}{8.068}{2.608}{%
\RG{0.000}{-0.003}{10.9091}{mi}{21}\RG{6.364}{-1.639}{7.9701}{rm}{49}\RG{11.096}{-0.003}{10.9091}{mi}{59}\RG{15.929}{-0.003}{10.9091}{mi}{58}\RG{20.773}{-0.003}{10.9091}{mi}{58}\RG{25.605}{-0.003}{10.9091}{mi}{58}\RG{30.449}{-0.003}{10.9091}{mi}{59}\RG{35.293}{-0.003}{10.9091}{mi}{21}\RG{41.702}{-1.639}{7.9701}{mi}{114}%
}\hskip.1em\allowbreak{} 两两不同。定义插值多项式

\[
\RB{171.786}{15.446}{15.572}{%
\RG{39.712}{10.359}{10.9092}{ex}{89}\RG{59.498}{7.375}{10.9091}{mi}{120}\RG{68.157}{7.375}{10.9091}{sy}{0}\RG{79.066}{7.375}{10.9091}{mi}{21}\RG{85.430}{5.598}{7.9701}{mi}{107}\RG{0.000}{-0.005}{10.9091}{mi}{112}\RG{9.871}{-0.005}{10.9091}{rm}{40}\RG{14.113}{-0.005}{10.9091}{mi}{120}\RG{20.348}{-0.005}{10.9091}{rm}{41}\RG{27.603}{-0.005}{10.9091}{rm}{61}\RG{115.819}{-0.005}{10.9091}{rm}{40}\RG{120.052}{-0.005}{10.9091}{rm}{49}\RG{128.546}{-0.005}{10.9091}{sy}{20}\RG{140.061}{-0.005}{10.9091}{mi}{106}\RG{148.208}{-0.005}{10.9091}{sy}{20}\RG{159.723}{-0.005}{10.9091}{mi}{114}\RG{164.948}{-0.005}{10.9091}{rm}{41}\RG{169.191}{-0.005}{10.9091}{mi}{58}\RG{5.489}{-1.642}{7.9701}{mi}{106}\RG{57.242}{-7.488}{10.9091}{mi}{21}\RG{70.412}{-7.488}{10.9091}{sy}{0}\RG{81.322}{-7.488}{10.9091}{mi}{21}\RG{63.605}{-9.125}{7.9701}{mi}{106}\RG{87.685}{-9.265}{7.9701}{mi}{107}\RG{39.136}{-13.358}{7.9701}{mi}{107}\RG{43.758}{-13.358}{7.9701}{sy}{54}\RG{43.758}{-13.358}{7.9701}{rm}{61}\RG{50.344}{-13.358}{7.9701}{mi}{106}\RR{57.242}{2.722}{35.562}{0.436}%
}
\]
由于\allowbreak{}\hskip.1em \RB{60.930}{8.678}{5.356}{%
\RG{0.000}{8.178}{10.9092}{ex}{80}\RG{11.515}{-3.230}{7.9701}{mi}{106}\RG{17.716}{-0.004}{10.9091}{mi}{112}\RG{23.204}{-1.640}{7.9701}{mi}{106}\RG{27.587}{-0.004}{10.9091}{rm}{40}\RG{31.829}{-0.004}{10.9091}{mi}{120}\RG{38.064}{-0.004}{10.9091}{rm}{41}\RG{44.840}{-0.004}{10.9091}{sy}{0}\RG{55.859}{-0.004}{10.9091}{rm}{49}%
}\hskip.1em\allowbreak{} 的次数不超过\allowbreak{}\hskip.1em \RB{23.848}{7.762}{0.624}{%
\RG{0.000}{-0.004}{10.9091}{mi}{114}\RG{7.759}{-0.004}{10.9091}{sy}{0}\RG{18.777}{-0.004}{10.9091}{rm}{49}%
}\hskip.1em\allowbreak{} ，且在\allowbreak{}\hskip.1em \RB{5.256}{5.318}{0.624}{%
\RG{0.000}{-0.004}{10.9091}{mi}{114}%
}\hskip.1em\allowbreak{} 个不同点\allowbreak{}\hskip.1em \RB{45.885}{8.067}{2.609}{%
\RG{0.000}{-0.004}{10.9091}{mi}{21}\RG{6.364}{-1.640}{7.9701}{rm}{49}\RG{11.096}{-0.004}{10.9091}{mi}{59}\RG{15.929}{-0.004}{10.9091}{mi}{58}\RG{20.773}{-0.004}{10.9091}{mi}{58}\RG{25.616}{-0.004}{10.9091}{mi}{58}\RG{30.449}{-0.004}{10.9091}{mi}{59}\RG{35.293}{-0.004}{10.9091}{mi}{21}\RG{41.702}{-1.640}{7.9701}{mi}{114}%
}\hskip.1em\allowbreak{} 处均为零，故\allowbreak{}\hskip.1em \RB{62.444}{8.678}{5.356}{%
\RG{0.000}{8.178}{10.9092}{ex}{80}\RG{11.515}{-3.230}{7.9701}{mi}{106}\RG{17.715}{-0.004}{10.9091}{mi}{112}\RG{23.204}{-1.640}{7.9701}{mi}{106}\RG{27.586}{-0.004}{10.9091}{rm}{40}\RG{31.828}{-0.004}{10.9091}{mi}{120}\RG{38.063}{-0.004}{10.9091}{rm}{41}\RG{45.591}{-0.004}{10.9091}{rm}{61}\RG{57.373}{-0.004}{10.9091}{rm}{49}%
}\hskip.1em\allowbreak{} 。此外，

\[
\RB{159.401}{16.059}{12.925}{%
\RG{106.869}{7.377}{10.9091}{mi}{104}\RG{113.154}{7.377}{10.9091}{rm}{40}\RG{117.396}{7.377}{10.9091}{mi}{120}\RG{123.631}{7.377}{10.9091}{rm}{41}\RG{85.670}{0.696}{10.9092}{ex}{81}\RG{0.000}{-0.003}{10.9091}{rm}{40}\RG{4.242}{-0.003}{10.9091}{mi}{120}\RG{12.901}{-0.003}{10.9091}{sy}{0}\RG{23.811}{-0.003}{10.9091}{mi}{21}\RG{34.556}{-0.003}{10.9091}{rm}{41}\RG{38.799}{-0.003}{10.9091}{mi}{112}\RG{48.670}{-0.003}{10.9091}{rm}{40}\RG{52.912}{-0.003}{10.9091}{mi}{120}\RG{59.148}{-0.003}{10.9091}{rm}{41}\RG{66.402}{-0.003}{10.9091}{rm}{61}\RG{156.806}{-0.003}{10.9091}{mi}{58}\RG{30.174}{-1.639}{7.9701}{mi}{106}\RG{44.287}{-1.639}{7.9701}{mi}{106}\RG{79.131}{-7.486}{10.9091}{mi}{99}\RG{111.564}{-7.486}{10.9091}{rm}{40}\RG{115.806}{-7.486}{10.9091}{mi}{21}\RG{128.976}{-7.486}{10.9091}{sy}{0}\RG{139.885}{-7.486}{10.9091}{mi}{21}\RG{151.369}{-7.486}{10.9091}{rm}{41}\RG{122.170}{-9.122}{7.9701}{mi}{106}\RG{146.249}{-9.263}{7.9701}{mi}{107}\RG{95.973}{-10.711}{7.9701}{mi}{107}\RG{100.595}{-10.711}{7.9701}{sy}{54}\RG{100.595}{-10.711}{7.9701}{rm}{61}\RG{107.181}{-10.711}{7.9701}{mi}{106}\RR{79.131}{2.725}{76.481}{0.436}%
}
\]
将\allowbreak{}\hskip.1em \RB{6.249}{5.326}{0.616}{%
\RG{0.000}{0.004}{10.9091}{mi}{120}%
}\hskip.1em\allowbreak{} 换为\allowbreak{}\hskip.1em \RB{8.366}{8.315}{0.500}{%
\RG{0.000}{0.004}{10.9091}{mi}{65}%
}\hskip.1em\allowbreak{} ，得到

\[
\RB{196.068}{17.655}{15.209}{%
\RG{30.827}{13.640}{7.9701}{mi}{114}\RG{24.977}{10.367}{10.9092}{ex}{88}\RG{0.000}{0.004}{10.9091}{mi}{73}\RG{13.462}{0.004}{10.9091}{rm}{61}\RG{42.553}{0.004}{10.9091}{mi}{112}\RG{52.423}{0.004}{10.9091}{rm}{40}\RG{56.666}{0.004}{10.9091}{mi}{65}\RG{64.849}{0.004}{10.9091}{rm}{41}\RG{69.091}{0.004}{10.9091}{mi}{59}\RG{95.758}{0.004}{10.9091}{rm}{40}\RG{100.000}{0.004}{10.9091}{mi}{65}\RG{110.606}{0.004}{10.9091}{sy}{0}\RG{121.515}{0.004}{10.9091}{mi}{21}\RG{132.261}{0.004}{10.9091}{mi}{73}\RG{142.693}{0.004}{10.9091}{rm}{41}\RG{146.935}{0.004}{10.9091}{mi}{112}\RG{156.806}{0.004}{10.9091}{rm}{40}\RG{161.049}{0.004}{10.9091}{mi}{65}\RG{169.230}{0.004}{10.9091}{rm}{41}\RG{176.485}{0.004}{10.9091}{rm}{61}\RG{187.994}{0.004}{10.9091}{rm}{48}\RG{193.474}{0.004}{10.9091}{mi}{58}\RG{4.795}{-1.633}{7.9701}{mi}{110}\RG{48.041}{-1.633}{7.9701}{mi}{106}\RG{127.879}{-1.633}{7.9701}{mi}{106}\RG{137.057}{-1.633}{7.9701}{mi}{110}\RG{152.424}{-1.633}{7.9701}{mi}{106}\RG{25.503}{-13.083}{7.9701}{mi}{106}\RG{29.387}{-13.083}{7.9701}{rm}{61}\RG{35.971}{-13.083}{7.9701}{rm}{49}%
}
\]
因此每个向量\allowbreak{}\hskip.1em \RB{5.595}{5.320}{0.622}{%
\RG{0.000}{-0.002}{10.9091}{mi}{122}%
}\hskip.1em\allowbreak{} 都可写为\allowbreak{}\hskip.1em \RB{70.326}{8.680}{5.353}{%
\RG{0.000}{-0.002}{10.9091}{mi}{122}\RG{8.773}{-0.002}{10.9091}{rm}{61}\RG{20.478}{8.180}{10.9092}{ex}{80}\RG{31.993}{-3.227}{7.9701}{mi}{106}\RG{38.194}{-0.002}{10.9091}{mi}{112}\RG{43.682}{-1.638}{7.9701}{mi}{106}\RG{48.065}{-0.002}{10.9091}{rm}{40}\RG{52.307}{-0.002}{10.9091}{mi}{65}\RG{60.490}{-0.002}{10.9091}{rm}{41}\RG{64.732}{-0.002}{10.9091}{mi}{122}%
}\hskip.1em\allowbreak{} ，其中每个非零项都是\allowbreak{}\hskip.1em \RB{8.366}{8.309}{0.502}{%
\RG{0.000}{-0.002}{10.9091}{mi}{65}%
}\hskip.1em\allowbreak{} 的特征向量。故特征向量张成\allowbreak{}\hskip.1em \RB{13.257}{8.183}{0.704}{%
\RG{0.000}{0.003}{10.9091}{bb}{67}\RG{7.879}{3.962}{7.9701}{mi}{110}%
}\hskip.1em\allowbreak{} ，从中可取一组基，\allowbreak{}\hskip.1em \RB{8.366}{8.314}{0.500}{%
\RG{0.000}{0.003}{10.9091}{mi}{65}%
}\hskip.1em\allowbreak{} 可对角化。\allowbreak{}\hskip.1em \RB{8.366}{7.998}{0.500}{%
\RG{0.000}{0.003}{10.9091}{am}{3}%
}\hskip.1em\allowbreak{}

小结：秩条件给出两个核空间的直和分解，互素与无重根条件进一步给出完整的特征向量系。

题目来源：谢启鸿《复旦大学高等代数习题集（2023版）》，第93页。

% Original page 6.
\section*{2020S05}
设\allowbreak{}\hskip.1em \RB{6.729}{5.323}{0.619}{%
\RG{0.000}{0.001}{10.9091}{mi}{110}%
}\hskip.1em\allowbreak{} 阶实方阵\allowbreak{}\hskip.1em \RB{21.762}{8.312}{2.604}{%
\RG{0.000}{0.001}{10.9091}{mi}{65}\RG{8.182}{0.001}{10.9091}{mi}{59}\RG{13.015}{0.001}{10.9091}{mi}{66}%
}\hskip.1em\allowbreak{} 的特征值都大于零，且

\[
\RB{124.658}{10.382}{1.408}{%
\RG{8.182}{4.502}{7.9701}{rm}{52}\RG{39.884}{4.502}{7.9701}{rm}{51}\RG{90.443}{4.502}{7.9701}{rm}{51}\RG{117.330}{4.502}{7.9701}{rm}{52}\RG{0.000}{-0.003}{10.9091}{mi}{65}\RG{15.338}{-0.003}{10.9091}{rm}{43}\RG{26.236}{-0.003}{10.9091}{rm}{50}\RG{31.702}{-0.003}{10.9091}{mi}{65}\RG{44.616}{-0.003}{10.9091}{mi}{66}\RG{56.468}{-0.003}{10.9091}{rm}{61}\RG{67.967}{-0.003}{10.9091}{rm}{50}\RG{73.439}{-0.003}{10.9091}{mi}{65}\RG{81.621}{-0.003}{10.9091}{mi}{66}\RG{97.599}{-0.003}{10.9091}{rm}{43}\RG{108.508}{-0.003}{10.9091}{mi}{66}\RG{122.063}{-0.003}{10.9091}{mi}{58}%
}
\]
证明：\allowbreak{}\hskip.1em \RB{31.475}{8.314}{0.500}{%
\RG{0.000}{0.003}{10.9091}{mi}{65}\RG{11.212}{0.003}{10.9091}{rm}{61}\RG{22.727}{0.003}{10.9091}{mi}{66}%
}\hskip.1em\allowbreak{} 。

证明：先证明线性映射

\[
\RB{224.955}{8.680}{3.229}{%
\RG{0.000}{-0.002}{10.9091}{sy}{84}\RG{11.744}{-0.002}{10.9091}{rm}{58}\RG{17.805}{-0.002}{10.9091}{mi}{77}\RG{34.025}{-0.002}{10.9091}{rm}{40}\RG{38.267}{-0.002}{10.9091}{bb}{67}\RG{46.146}{-0.002}{10.9091}{rm}{41}\RG{53.418}{-0.002}{10.9091}{sy}{0}\RG{60.084}{-0.002}{10.9091}{sy}{33}\RG{74.026}{-0.002}{10.9091}{mi}{77}\RG{90.245}{-0.002}{10.9091}{rm}{40}\RG{94.488}{-0.002}{10.9091}{bb}{67}\RG{102.367}{-0.002}{10.9091}{rm}{41}\RG{106.609}{-0.002}{10.9091}{mi}{59}\RG{133.276}{-0.002}{10.9091}{sy}{84}\RG{141.990}{-0.002}{10.9091}{rm}{40}\RG{146.232}{-0.002}{10.9091}{mi}{88}\RG{156.126}{-0.002}{10.9091}{rm}{41}\RG{163.381}{-0.002}{10.9091}{rm}{61}\RG{174.914}{-0.002}{10.9091}{mi}{65}\RG{183.096}{-0.002}{10.9091}{mi}{88}\RG{195.415}{-0.002}{10.9091}{rm}{43}\RG{206.324}{-0.002}{10.9091}{mi}{88}\RG{216.207}{-0.002}{10.9091}{mi}{66}\RG{28.388}{-1.638}{7.9701}{mi}{110}\RG{84.609}{-1.638}{7.9701}{mi}{110}%
}
\]
的核为零。

设\allowbreak{}\hskip.1em \RB{70.172}{8.311}{1.405}{%
\RG{0.000}{-0.000}{10.9091}{mi}{65}\RG{8.182}{-0.000}{10.9091}{mi}{88}\RG{20.500}{-0.000}{10.9091}{rm}{43}\RG{31.409}{-0.000}{10.9091}{mi}{88}\RG{41.293}{-0.000}{10.9091}{mi}{66}\RG{53.155}{-0.000}{10.9091}{rm}{61}\RG{64.654}{-0.000}{10.9091}{rm}{48}%
}\hskip.1em\allowbreak{} 。将\allowbreak{}\hskip.1em \RB{8.747}{7.951}{0.500}{%
\RG{0.000}{-0.000}{10.9091}{mi}{66}%
}\hskip.1em\allowbreak{} 相似化为上三角矩阵

\[
\RB{204.394}{10.301}{3.766}{%
\RG{48.506}{4.501}{7.9701}{sy}{0}\RG{55.092}{4.501}{7.9701}{rm}{49}\RG{0.000}{-0.004}{10.9091}{mi}{66}\RG{11.852}{-0.004}{10.9091}{rm}{61}\RG{23.367}{-0.004}{10.9091}{mi}{81}\RG{31.985}{-0.004}{10.9091}{mi}{84}\RG{39.862}{-0.004}{10.9091}{mi}{81}\RG{59.826}{-0.004}{10.9091}{mi}{59}\RG{86.477}{-0.004}{10.9091}{mi}{84}\RG{97.413}{-0.004}{10.9091}{rm}{61}\RG{108.911}{-0.004}{10.9091}{rm}{40}\RG{113.170}{-0.004}{10.9091}{mi}{116}\RG{124.375}{-0.004}{10.9091}{rm}{41}\RG{128.618}{-0.004}{10.9091}{mi}{59}\RG{144.360}{-0.004}{10.9091}{mi}{116}\RG{159.611}{-0.004}{10.9091}{rm}{61}\RG{171.126}{-0.004}{10.9091}{mi}{12}\RG{184.710}{-0.004}{10.9091}{mi}{62}\RG{196.225}{-0.004}{10.9091}{rm}{48}\RG{201.680}{-0.004}{10.9091}{mi}{59}\RG{117.110}{-1.640}{7.9701}{mi}{105}\RG{119.987}{-1.640}{7.9701}{mi}{106}\RG{148.314}{-1.640}{7.9701}{mi}{106}\RG{152.188}{-1.640}{7.9701}{mi}{106}\RG{177.297}{-1.640}{7.9701}{mi}{106}%
}
\]
并令\allowbreak{}\hskip.1em \RB{109.571}{8.684}{3.225}{%
\RG{0.000}{0.002}{10.9091}{mi}{89}\RG{11.788}{0.002}{10.9091}{rm}{61}\RG{23.303}{0.002}{10.9091}{mi}{88}\RG{33.187}{0.002}{10.9091}{mi}{81}\RG{44.851}{0.002}{10.9091}{rm}{61}\RG{56.350}{0.002}{10.9091}{rm}{40}\RG{60.609}{0.002}{10.9091}{mi}{121}\RG{65.958}{-1.634}{7.9701}{rm}{49}\RG{70.691}{0.002}{10.9091}{mi}{59}\RG{75.523}{0.002}{10.9091}{mi}{58}\RG{80.367}{0.002}{10.9091}{mi}{58}\RG{85.200}{0.002}{10.9091}{mi}{58}\RG{90.043}{0.002}{10.9091}{mi}{59}\RG{94.887}{0.002}{10.9091}{mi}{121}\RG{100.282}{-1.634}{7.9701}{mi}{110}\RG{105.918}{0.002}{10.9091}{rm}{41}%
}\hskip.1em\allowbreak{} 。由\allowbreak{}\hskip.1em \RB{66.968}{8.313}{1.403}{%
\RG{0.000}{0.002}{10.9091}{mi}{65}\RG{8.182}{0.002}{10.9091}{mi}{89}\RG{19.364}{0.002}{10.9091}{rm}{43}\RG{30.273}{0.002}{10.9091}{mi}{89}\RG{39.022}{0.002}{10.9091}{mi}{84}\RG{49.951}{0.002}{10.9091}{rm}{61}\RG{61.450}{0.002}{10.9091}{rm}{48}%
}\hskip.1em\allowbreak{} ，比较第\allowbreak{}\hskip.1em \RB{4.831}{7.713}{2.734}{%
\RG{0.000}{0.002}{10.9091}{mi}{106}%
}\hskip.1em\allowbreak{} 列，得

\[
\RB{129.147}{10.862}{15.481}{%
\RG{85.565}{10.362}{10.9092}{ex}{88}\RG{0.000}{-0.002}{10.9091}{rm}{40}\RG{4.243}{-0.002}{10.9091}{mi}{65}\RG{14.849}{-0.002}{10.9091}{rm}{43}\RG{25.758}{-0.002}{10.9091}{mi}{12}\RG{36.311}{-0.002}{10.9091}{mi}{73}\RG{46.742}{-0.002}{10.9091}{rm}{41}\RG{50.985}{-0.002}{10.9091}{mi}{121}\RG{63.747}{-0.002}{10.9091}{rm}{61}\RG{75.262}{-0.002}{10.9091}{sy}{0}\RG{103.141}{-0.002}{10.9091}{mi}{116}\RG{116.084}{-0.002}{10.9091}{mi}{121}\RG{126.552}{-0.002}{10.9091}{mi}{58}\RG{31.928}{-1.638}{7.9701}{mi}{106}\RG{41.106}{-1.638}{7.9701}{mi}{110}\RG{56.333}{-1.638}{7.9701}{mi}{106}\RG{107.080}{-1.779}{7.9701}{mi}{107}\RG{111.695}{-1.779}{7.9701}{mi}{106}\RG{121.432}{-1.779}{7.9701}{mi}{107}\RG{85.898}{-13.355}{7.9701}{mi}{107}\RG{90.512}{-13.355}{7.9701}{mi}{60}\RG{97.096}{-13.355}{7.9701}{mi}{106}%
}
\]
\allowbreak{}\hskip.1em \RB{41.999}{8.308}{3.765}{%
\RG{0.000}{-0.003}{10.9091}{mi}{65}\RG{10.485}{-0.003}{10.9091}{rm}{43}\RG{21.273}{-0.003}{10.9091}{mi}{12}\RG{27.443}{-1.639}{7.9701}{mi}{106}\RG{31.825}{-0.003}{10.9091}{mi}{73}\RG{36.621}{-1.639}{7.9701}{mi}{110}%
}\hskip.1em\allowbreak{} 的特征值为\allowbreak{}\hskip.1em \RB{33.413}{8.199}{3.765}{%
\RG{0.000}{-0.003}{10.9091}{mi}{11}\RG{6.978}{-1.639}{7.9701}{mi}{105}\RG{12.663}{-0.003}{10.9091}{rm}{43}\RG{23.450}{-0.003}{10.9091}{mi}{12}\RG{29.621}{-1.639}{7.9701}{mi}{106}%
}\hskip.1em\allowbreak{} ，其中\allowbreak{}\hskip.1em \RB{30.423}{7.763}{2.219}{%
\RG{0.000}{-0.003}{10.9091}{mi}{11}\RG{6.978}{-1.639}{7.9701}{mi}{105}\RG{13.390}{-0.003}{10.9091}{mi}{62}\RG{24.905}{-0.003}{10.9091}{rm}{48}%
}\hskip.1em\allowbreak{} 是\allowbreak{}\hskip.1em \RB{8.366}{8.308}{0.503}{%
\RG{0.000}{-0.003}{10.9091}{mi}{65}%
}\hskip.1em\allowbreak{} 的特征值，故这些特征值均非零。于是\allowbreak{}\hskip.1em \RB{41.999}{8.308}{3.765}{%
\RG{0.000}{-0.003}{10.9091}{mi}{65}\RG{10.485}{-0.003}{10.9091}{rm}{43}\RG{21.272}{-0.003}{10.9091}{mi}{12}\RG{27.443}{-1.639}{7.9701}{mi}{106}\RG{31.825}{-0.003}{10.9091}{mi}{73}\RG{36.621}{-1.639}{7.9701}{mi}{110}%
}\hskip.1em\allowbreak{}可逆。

当\allowbreak{}\hskip.1em \RB{24.716}{7.764}{2.738}{%
\RG{0.000}{-0.002}{10.9091}{mi}{106}\RG{8.147}{-0.002}{10.9091}{rm}{61}\RG{19.645}{-0.002}{10.9091}{rm}{49}%
}\hskip.1em\allowbreak{} 时，上式右端为零，故\allowbreak{}\hskip.1em \RB{30.128}{7.764}{2.738}{%
\RG{0.000}{-0.002}{10.9091}{mi}{121}\RG{5.349}{-1.638}{7.9701}{rm}{49}\RG{13.111}{-0.002}{10.9091}{rm}{61}\RG{24.609}{-0.002}{10.9091}{rm}{48}%
}\hskip.1em\allowbreak{} 。若\allowbreak{}\hskip.1em \RB{92.498}{7.764}{3.764}{%
\RG{0.000}{-0.002}{10.9091}{mi}{121}\RG{5.349}{-1.638}{7.9701}{rm}{49}\RG{13.112}{-0.002}{10.9091}{rm}{61}\RG{24.627}{-0.002}{10.9091}{sy}{1}\RG{29.460}{-0.002}{10.9091}{sy}{1}\RG{34.304}{-0.002}{10.9091}{sy}{1}\RG{40.384}{-0.002}{10.9091}{rm}{61}\RG{51.899}{-0.002}{10.9091}{mi}{121}\RG{57.248}{-1.638}{7.9701}{mi}{106}\RG{61.133}{-1.638}{7.9701}{sy}{0}\RG{67.719}{-1.638}{7.9701}{rm}{49}\RG{75.482}{-0.002}{10.9091}{rm}{61}\RG{86.980}{-0.002}{10.9091}{rm}{48}%
}\hskip.1em\allowbreak{} ，则上式又给出\allowbreak{}\hskip.1em \RB{29.778}{7.764}{3.764}{%
\RG{0.000}{-0.002}{10.9091}{mi}{121}\RG{5.349}{-1.638}{7.9701}{mi}{106}\RG{12.761}{-0.002}{10.9091}{rm}{61}\RG{24.259}{-0.002}{10.9091}{rm}{48}%
}\hskip.1em\allowbreak{} 。归纳得\allowbreak{}\hskip.1em \RB{28.805}{7.954}{0.737}{%
\RG{0.000}{0.003}{10.9091}{mi}{89}\RG{11.788}{0.003}{10.9091}{rm}{61}\RG{23.286}{0.003}{10.9091}{rm}{48}%
}\hskip.1em\allowbreak{} ，从而\allowbreak{}\hskip.1em \RB{29.942}{7.954}{0.737}{%
\RG{0.000}{0.003}{10.9091}{mi}{88}\RG{12.925}{0.003}{10.9091}{rm}{61}\RG{24.423}{0.003}{10.9091}{rm}{48}%
}\hskip.1em\allowbreak{} 。因此\allowbreak{}\hskip.1em \RB{9.598}{8.325}{1.239}{%
\RG{0.000}{0.003}{10.9091}{sy}{84}%
}\hskip.1em\allowbreak{} 是单射。

现在令\allowbreak{}\hskip.1em \RB{58.579}{8.310}{2.304}{%
\RG{0.000}{-0.001}{10.9091}{mi}{88}\RG{9.038}{-1.637}{7.9701}{rm}{48}\RG{16.801}{-0.001}{10.9091}{rm}{61}\RG{28.316}{-0.001}{10.9091}{mi}{65}\RG{38.922}{-0.001}{10.9091}{sy}{0}\RG{49.831}{-0.001}{10.9091}{mi}{66}%
}\hskip.1em\allowbreak{} 。始终保留各矩阵因子的次序，依次计算：

\[
\RB{203.697}{8.677}{95.550}{%
\RG{0.000}{-0.005}{10.9091}{mi}{88}\RG{16.800}{-0.005}{10.9091}{rm}{58}\RG{19.822}{-0.005}{10.9091}{rm}{61}\RG{31.345}{-0.005}{10.9091}{sy}{84}\RG{40.060}{-0.005}{10.9091}{rm}{40}\RG{44.302}{-0.005}{10.9091}{mi}{88}\RG{58.073}{-0.005}{10.9091}{rm}{41}\RG{65.327}{-0.005}{10.9091}{rm}{61}\RG{76.861}{-0.005}{10.9091}{mi}{65}\RG{85.043}{-0.005}{10.9091}{rm}{40}\RG{89.285}{-0.005}{10.9091}{mi}{65}\RG{99.891}{-0.005}{10.9091}{sy}{0}\RG{110.801}{-0.005}{10.9091}{mi}{66}\RG{119.623}{-0.005}{10.9091}{rm}{41}\RG{126.278}{-0.005}{10.9091}{rm}{43}\RG{137.176}{-0.005}{10.9091}{rm}{40}\RG{141.441}{-0.005}{10.9091}{mi}{65}\RG{152.047}{-0.005}{10.9091}{sy}{0}\RG{162.956}{-0.005}{10.9091}{mi}{66}\RG{171.778}{-0.005}{10.9091}{rm}{41}\RG{176.020}{-0.005}{10.9091}{mi}{66}\RG{9.037}{-1.641}{7.9701}{rm}{49}\RG{53.340}{-1.641}{7.9701}{rm}{48}\RG{85.043}{-13.664}{7.9701}{rm}{50}\RG{111.931}{-13.664}{7.9701}{rm}{50}\RG{65.346}{-18.169}{10.9091}{rm}{61}\RG{76.861}{-18.169}{10.9091}{mi}{65}\RG{92.200}{-18.169}{10.9091}{sy}{0}\RG{103.109}{-18.169}{10.9091}{mi}{66}\RG{116.663}{-18.169}{10.9091}{mi}{59}\RG{97.467}{-33.821}{7.9701}{rm}{50}\RG{124.355}{-33.821}{7.9701}{rm}{50}\RG{159.087}{-33.821}{7.9701}{rm}{50}\RG{185.976}{-33.821}{7.9701}{rm}{50}\RG{0.000}{-38.325}{10.9091}{mi}{88}\RG{16.800}{-38.325}{10.9091}{rm}{58}\RG{19.822}{-38.325}{10.9091}{rm}{61}\RG{31.345}{-38.325}{10.9091}{sy}{84}\RG{40.060}{-38.325}{10.9091}{rm}{40}\RG{44.302}{-38.325}{10.9091}{mi}{88}\RG{58.072}{-38.325}{10.9091}{rm}{41}\RG{65.327}{-38.325}{10.9091}{rm}{61}\RG{76.861}{-38.325}{10.9091}{mi}{65}\RG{85.043}{-38.325}{10.9091}{rm}{40}\RG{89.285}{-38.325}{10.9091}{mi}{65}\RG{104.624}{-38.325}{10.9091}{sy}{0}\RG{115.533}{-38.325}{10.9091}{mi}{66}\RG{129.087}{-38.325}{10.9091}{rm}{41}\RG{135.742}{-38.325}{10.9091}{rm}{43}\RG{146.640}{-38.325}{10.9091}{rm}{40}\RG{150.905}{-38.325}{10.9091}{mi}{65}\RG{166.244}{-38.325}{10.9091}{sy}{0}\RG{177.153}{-38.325}{10.9091}{mi}{66}\RG{190.708}{-38.325}{10.9091}{rm}{41}\RG{194.950}{-38.325}{10.9091}{mi}{66}\RG{9.037}{-39.961}{7.9701}{rm}{50}\RG{53.340}{-39.961}{7.9701}{rm}{49}\RG{85.043}{-51.984}{7.9701}{rm}{51}\RG{111.290}{-51.984}{7.9701}{rm}{50}\RG{155.182}{-51.984}{7.9701}{rm}{50}\RG{182.069}{-51.984}{7.9701}{rm}{51}\RG{65.346}{-56.489}{10.9091}{rm}{61}\RG{76.861}{-56.489}{10.9091}{mi}{65}\RG{92.200}{-56.489}{10.9091}{rm}{43}\RG{103.109}{-56.489}{10.9091}{mi}{65}\RG{116.023}{-56.489}{10.9091}{mi}{66}\RG{127.269}{-56.489}{10.9091}{sy}{0}\RG{138.178}{-56.489}{10.9091}{mi}{65}\RG{146.360}{-56.489}{10.9091}{mi}{66}\RG{162.338}{-56.489}{10.9091}{sy}{0}\RG{173.247}{-56.489}{10.9091}{mi}{66}\RG{186.803}{-56.489}{10.9091}{mi}{59}\RG{85.043}{-72.142}{7.9701}{rm}{52}\RG{116.745}{-72.142}{7.9701}{rm}{51}\RG{166.091}{-72.142}{7.9701}{rm}{51}\RG{192.980}{-72.142}{7.9701}{rm}{52}\RG{31.345}{-76.646}{10.9091}{sy}{84}\RG{40.060}{-76.646}{10.9091}{rm}{40}\RG{44.302}{-76.646}{10.9091}{mi}{88}\RG{58.072}{-76.646}{10.9091}{rm}{41}\RG{65.327}{-76.646}{10.9091}{rm}{61}\RG{76.861}{-76.646}{10.9091}{mi}{65}\RG{92.200}{-76.646}{10.9091}{rm}{43}\RG{103.098}{-76.646}{10.9091}{rm}{50}\RG{108.563}{-76.646}{10.9091}{mi}{65}\RG{121.477}{-76.646}{10.9091}{mi}{66}\RG{132.724}{-76.646}{10.9091}{sy}{0}\RG{143.633}{-76.646}{10.9091}{rm}{50}\RG{149.087}{-76.646}{10.9091}{mi}{65}\RG{157.269}{-76.646}{10.9091}{mi}{66}\RG{173.248}{-76.646}{10.9091}{sy}{0}\RG{184.158}{-76.646}{10.9091}{mi}{66}\RG{53.340}{-78.282}{7.9701}{rm}{50}\RG{65.346}{-94.810}{10.9091}{rm}{61}\RG{76.844}{-94.810}{10.9091}{rm}{48}\RG{82.316}{-94.810}{10.9091}{mi}{58}%
}
\]
由\allowbreak{}\hskip.1em \RB{9.598}{8.319}{1.245}{%
\RG{0.000}{-0.003}{10.9091}{sy}{84}%
}\hskip.1em\allowbreak{} 的单射性，逐次得到

\[
\RB{157.240}{7.953}{2.604}{%
\RG{0.000}{0.002}{10.9091}{mi}{88}\RG{16.801}{0.002}{10.9091}{rm}{61}\RG{28.299}{0.002}{10.9091}{rm}{48}\RG{33.770}{0.002}{10.9091}{mi}{59}\RG{60.421}{0.002}{10.9091}{mi}{88}\RG{77.238}{0.002}{10.9091}{rm}{61}\RG{88.736}{0.002}{10.9091}{rm}{48}\RG{94.208}{0.002}{10.9091}{mi}{59}\RG{120.859}{0.002}{10.9091}{mi}{88}\RG{137.675}{0.002}{10.9091}{rm}{61}\RG{149.173}{0.002}{10.9091}{rm}{48}\RG{154.645}{0.002}{10.9091}{mi}{58}\RG{9.038}{-1.634}{7.9701}{rm}{50}\RG{69.476}{-1.634}{7.9701}{rm}{49}\RG{129.913}{-1.634}{7.9701}{rm}{48}%
}
\]
于是\allowbreak{}\hskip.1em \RB{50.384}{8.312}{0.739}{%
\RG{0.000}{0.001}{10.9091}{mi}{65}\RG{10.606}{0.001}{10.9091}{sy}{0}\RG{21.515}{0.001}{10.9091}{mi}{66}\RG{33.367}{0.001}{10.9091}{rm}{61}\RG{44.866}{0.001}{10.9091}{rm}{48}%
}\hskip.1em\allowbreak{} ，即\allowbreak{}\hskip.1em \RB{31.475}{8.312}{0.500}{%
\RG{0.000}{0.001}{10.9091}{mi}{65}\RG{11.212}{0.001}{10.9091}{rm}{61}\RG{22.727}{0.001}{10.9091}{mi}{66}%
}\hskip.1em\allowbreak{} 。\allowbreak{}\hskip.1em \RB{8.366}{7.996}{0.500}{%
\RG{0.000}{0.001}{10.9091}{am}{3}%
}\hskip.1em\allowbreak{}

小结：原等式就是\allowbreak{}\hskip.1em \RB{72.308}{9.759}{3.227}{%
\RG{0.000}{-0.000}{10.9091}{sy}{84}\RG{8.714}{3.959}{7.9701}{rm}{51}\RG{13.446}{-0.000}{10.9091}{rm}{40}\RG{17.689}{-0.000}{10.9091}{mi}{65}\RG{28.295}{-0.000}{10.9091}{sy}{0}\RG{39.204}{-0.000}{10.9091}{mi}{66}\RG{48.026}{-0.000}{10.9091}{rm}{41}\RG{55.281}{-0.000}{10.9091}{rm}{61}\RG{66.790}{-0.000}{10.9091}{rm}{48}%
}\hskip.1em\allowbreak{} ；正特征值条件保证\allowbreak{}\hskip.1em \RB{9.598}{8.322}{1.242}{%
\RG{0.000}{-0.000}{10.9091}{sy}{84}%
}\hskip.1em\allowbreak{} 的单射性，从而可以连续三次消去\allowbreak{}\hskip.1em \RB{9.598}{8.322}{1.242}{%
\RG{0.000}{-0.000}{10.9091}{sy}{84}%
}\hskip.1em\allowbreak{} 。

题目来源：谢启鸿《复旦大学高等代数习题集（2023版）》，第93页。

% Original page 7.
\section*{2020S06}
设\allowbreak{}\hskip.1em \RB{53.509}{8.683}{3.226}{%
\RG{0.000}{0.001}{10.9091}{mi}{65}\RG{11.212}{0.001}{10.9091}{sy}{50}\RG{21.515}{0.001}{10.9091}{mi}{77}\RG{32.098}{-1.635}{7.9701}{mi}{110}\RG{37.735}{0.001}{10.9091}{rm}{40}\RG{41.977}{0.001}{10.9091}{bb}{67}\RG{49.856}{0.001}{10.9091}{rm}{41}%
}\hskip.1em\allowbreak{} ，\allowbreak{}\hskip.1em \RB{53.947}{8.683}{3.226}{%
\RG{0.000}{0.001}{10.9091}{mi}{102}\RG{6.515}{0.001}{10.9091}{rm}{40}\RG{10.758}{0.001}{10.9091}{mi}{120}\RG{16.993}{0.001}{10.9091}{rm}{41}\RG{24.265}{0.001}{10.9091}{sy}{50}\RG{34.568}{0.001}{10.9091}{bb}{67}\RG{42.447}{0.001}{10.9091}{rm}{91}\RG{45.478}{0.001}{10.9091}{mi}{120}\RG{51.712}{0.001}{10.9091}{rm}{93}%
}\hskip.1em\allowbreak{} 。证明：\allowbreak{}\hskip.1em \RB{8.366}{8.312}{0.500}{%
\RG{0.000}{0.001}{10.9091}{mi}{65}%
}\hskip.1em\allowbreak{} 可对角化，当且仅当

\[
\RB{97.836}{12.309}{21.407}{%
\RG{26.318}{11.373}{10.9092}{ex}{32}\RG{91.282}{11.373}{10.9092}{ex}{33}\RG{42.455}{-0.002}{10.9091}{mi}{65}\RG{68.092}{-0.002}{10.9091}{mi}{102}\RG{74.615}{-0.002}{10.9091}{rm}{40}\RG{78.858}{-0.002}{10.9091}{mi}{65}\RG{87.040}{-0.002}{10.9091}{rm}{41}\RG{0.000}{-7.281}{10.9091}{mi}{77}\RG{14.803}{-7.281}{10.9091}{rm}{61}\RG{34.955}{-15.177}{10.9091}{mi}{102}\RG{41.470}{-15.177}{10.9091}{rm}{40}\RG{45.712}{-15.177}{10.9091}{mi}{65}\RG{53.894}{-15.177}{10.9091}{rm}{41}\RG{75.600}{-15.177}{10.9091}{mi}{65}%
}
\]
可对角化。

证明：令

\[
\RB{160.256}{19.589}{14.128}{%
\RG{23.065}{18.652}{10.9092}{ex}{32}\RG{71.013}{18.652}{10.9092}{ex}{33}\RG{143.713}{7.377}{10.9091}{rm}{49}\RG{31.701}{7.277}{10.9091}{mi}{73}\RG{56.338}{7.277}{10.9091}{mi}{73}\RG{36.497}{5.641}{7.9701}{mi}{110}\RG{61.135}{5.641}{7.9701}{mi}{110}\RG{116.653}{4.501}{7.9701}{sy}{0}\RG{123.240}{4.501}{7.9701}{rm}{49}\RG{0.000}{-0.003}{10.9091}{mi}{80}\RG{11.550}{-0.003}{10.9091}{rm}{61}\RG{81.468}{-0.003}{10.9091}{mi}{59}\RG{108.119}{-0.003}{10.9091}{mi}{80}\RG{131.003}{-0.003}{10.9091}{rm}{61}\RG{150.363}{-0.003}{10.9091}{mi}{80}\RG{157.661}{-0.003}{10.9091}{mi}{58}\RG{143.713}{-7.486}{10.9091}{rm}{50}\RG{31.701}{-7.898}{10.9091}{mi}{73}\RG{52.096}{-7.898}{10.9091}{sy}{0}\RG{60.581}{-7.898}{10.9091}{mi}{73}\RG{36.497}{-9.535}{7.9701}{mi}{110}\RG{65.377}{-9.535}{7.9701}{mi}{110}\RR{143.713}{2.724}{5.455}{0.436}%
}
\]
直接作分块乘法，得到

\[
\RB{175.719}{12.312}{21.404}{%
\RG{54.675}{11.376}{10.9092}{ex}{32}\RG{162.669}{11.376}{10.9092}{ex}{33}\RG{63.313}{0.001}{10.9091}{mi}{65}\RG{73.919}{0.001}{10.9091}{rm}{43}\RG{84.828}{0.001}{10.9091}{mi}{102}\RG{91.343}{0.001}{10.9091}{rm}{40}\RG{95.585}{0.001}{10.9091}{mi}{65}\RG{103.767}{0.001}{10.9091}{rm}{41}\RG{137.574}{0.001}{10.9091}{rm}{48}\RG{8.519}{-2.775}{7.9701}{sy}{0}\RG{15.106}{-2.775}{7.9701}{rm}{49}\RG{0.000}{-7.279}{10.9091}{mi}{80}\RG{19.838}{-7.279}{10.9091}{mi}{77}\RG{31.609}{-7.279}{10.9091}{mi}{80}\RG{43.160}{-7.279}{10.9091}{rm}{61}\RG{173.124}{-7.279}{10.9091}{mi}{58}\RG{82.934}{-15.174}{10.9091}{rm}{48}\RG{117.972}{-15.174}{10.9091}{mi}{65}\RG{128.578}{-15.174}{10.9091}{sy}{0}\RG{139.487}{-15.174}{10.9091}{mi}{102}\RG{146.003}{-15.174}{10.9091}{rm}{40}\RG{150.245}{-15.174}{10.9091}{mi}{65}\RG{158.427}{-15.174}{10.9091}{rm}{41}%
}
\]
先证必要性。若\allowbreak{}\hskip.1em \RB{8.366}{8.315}{0.500}{%
\RG{0.000}{0.004}{10.9091}{mi}{65}%
}\hskip.1em\allowbreak{} 可对角化，取可逆矩阵\allowbreak{}\hskip.1em \RB{8.573}{8.195}{2.612}{%
\RG{0.000}{0.004}{10.9091}{mi}{81}%
}\hskip.1em\allowbreak{} ，使\allowbreak{}\hskip.1em \RB{60.565}{9.763}{2.612}{%
\RG{0.000}{0.004}{10.9091}{mi}{81}\RG{8.624}{3.963}{7.9701}{sy}{0}\RG{15.211}{3.963}{7.9701}{rm}{49}\RG{19.943}{0.004}{10.9091}{mi}{65}\RG{28.125}{0.004}{10.9091}{mi}{81}\RG{39.779}{0.004}{10.9091}{rm}{61}\RG{51.294}{0.004}{10.9091}{mi}{68}%
}\hskip.1em\allowbreak{} 为对角矩阵，则

\[
\RB{144.224}{10.305}{4.307}{%
\RG{19.943}{8.837}{10.9092}{ex}{0}\RG{69.641}{8.837}{10.9092}{ex}{1}\RG{8.624}{4.505}{7.9701}{sy}{0}\RG{15.210}{4.505}{7.9701}{rm}{49}\RG{0.000}{0.001}{10.9091}{mi}{81}\RG{24.943}{0.001}{10.9091}{mi}{65}\RG{35.549}{0.001}{10.9091}{sy}{6}\RG{46.458}{0.001}{10.9091}{mi}{102}\RG{52.973}{0.001}{10.9091}{rm}{40}\RG{57.215}{0.001}{10.9091}{mi}{65}\RG{65.398}{0.001}{10.9091}{rm}{41}\RG{74.641}{0.001}{10.9091}{mi}{81}\RG{86.295}{0.001}{10.9091}{rm}{61}\RG{97.810}{0.001}{10.9091}{mi}{68}\RG{109.569}{0.001}{10.9091}{sy}{6}\RG{120.478}{0.001}{10.9091}{mi}{102}\RG{126.994}{0.001}{10.9091}{rm}{40}\RG{131.236}{0.001}{10.9091}{mi}{68}\RG{140.571}{0.001}{10.9091}{rm}{41}%
}
\]
也是对角矩阵。因此用diag\allowbreak{}\hskip.1em \RB{29.992}{8.680}{3.229}{%
\RG{0.000}{-0.002}{10.9091}{rm}{40}\RG{4.242}{-0.002}{10.9091}{mi}{81}\RG{12.860}{-0.002}{10.9091}{mi}{59}\RG{17.693}{-0.002}{10.9091}{mi}{81}\RG{26.339}{-0.002}{10.9091}{rm}{41}%
}\hskip.1em\allowbreak{} 对上式右侧作相似变换，便得到对角矩阵，故\allowbreak{}\hskip.1em \RB{11.889}{7.949}{0.502}{%
\RG{0.000}{-0.002}{10.9091}{mi}{77}%
}\hskip.1em\allowbreak{} 可对角化。

再证充分性。先说明一个事实：可对角化的线性变换在任一不变子空间上的限制仍可对角化。设该变换为\allowbreak{}\hskip.1em \RB{8.180}{7.885}{0.501}{%
\RG{0.000}{-0.001}{10.9091}{mi}{84}%
}\hskip.1em\allowbreak{} ，其不同特征值为\allowbreak{}\hskip.1em \RB{45.885}{8.070}{2.606}{%
\RG{0.000}{-0.001}{10.9091}{mi}{21}\RG{6.363}{-1.637}{7.9701}{rm}{49}\RG{11.096}{-0.001}{10.9091}{mi}{59}\RG{15.929}{-0.001}{10.9091}{mi}{58}\RG{20.773}{-0.001}{10.9091}{mi}{58}\RG{25.605}{-0.001}{10.9091}{mi}{58}\RG{30.449}{-0.001}{10.9091}{mi}{59}\RG{35.293}{-0.001}{10.9091}{mi}{21}\RG{41.702}{-1.637}{7.9701}{mi}{114}%
}\hskip.1em\allowbreak{} 。投影到相应特征子空间的映射为

\[
\RB{87.235}{8.069}{22.949}{%
\RG{27.557}{2.981}{10.9092}{ex}{89}\RG{45.086}{-0.002}{10.9091}{mi}{84}\RG{55.400}{-0.002}{10.9091}{sy}{0}\RG{66.310}{-0.002}{10.9091}{mi}{21}\RG{77.794}{-0.002}{10.9091}{mi}{73}\RG{72.674}{-1.779}{7.9701}{mi}{107}\RG{0.000}{-7.382}{10.9091}{mi}{69}\RG{15.465}{-7.382}{10.9091}{rm}{61}\RG{84.641}{-7.382}{10.9091}{mi}{58}\RG{8.053}{-9.019}{7.9701}{mi}{106}\RG{46.484}{-14.865}{10.9091}{mi}{21}\RG{59.655}{-14.865}{10.9091}{sy}{0}\RG{70.564}{-14.865}{10.9091}{mi}{21}\RG{52.848}{-16.502}{7.9701}{mi}{106}\RG{76.928}{-16.643}{7.9701}{mi}{107}\RG{26.980}{-20.735}{7.9701}{mi}{107}\RG{31.602}{-20.735}{7.9701}{sy}{54}\RG{31.602}{-20.735}{7.9701}{rm}{61}\RG{38.188}{-20.735}{7.9701}{mi}{106}\RR{45.086}{-4.655}{38.358}{0.436}%
}
\]
若\allowbreak{}\hskip.1em \RB{8.791}{7.952}{0.739}{%
\RG{0.000}{0.001}{10.9091}{mi}{85}%
}\hskip.1em\allowbreak{} 是\allowbreak{}\hskip.1em \RB{8.180}{7.886}{0.500}{%
\RG{0.000}{0.001}{10.9091}{mi}{84}%
}\hskip.1em\allowbreak{} 的不变子空间，则\allowbreak{}\hskip.1em \RB{44.410}{7.952}{3.761}{%
\RG{0.000}{0.001}{10.9091}{mi}{69}\RG{8.054}{-1.635}{7.9701}{mi}{106}\RG{12.436}{0.001}{10.9091}{mi}{85}\RG{24.104}{0.001}{10.9091}{sy}{18}\RG{35.619}{0.001}{10.9091}{mi}{85}%
}\hskip.1em\allowbreak{} ，因为\allowbreak{}\hskip.1em \RB{11.845}{7.919}{3.761}{%
\RG{0.000}{0.001}{10.9091}{mi}{69}\RG{8.053}{-1.635}{7.9701}{mi}{106}%
}\hskip.1em\allowbreak{} 是\allowbreak{}\hskip.1em \RB{8.180}{7.886}{0.500}{%
\RG{0.000}{0.001}{10.9091}{mi}{84}%
}\hskip.1em\allowbreak{} 的多项式。由\allowbreak{}\hskip.1em \RB{50.618}{8.683}{5.350}{%
\RG{0.000}{8.183}{10.9092}{ex}{80}\RG{11.515}{-3.224}{7.9701}{mi}{106}\RG{17.715}{0.001}{10.9091}{mi}{69}\RG{25.768}{-1.635}{7.9701}{mi}{106}\RG{33.181}{0.001}{10.9091}{rm}{61}\RG{44.696}{0.001}{10.9091}{mi}{73}%
}\hskip.1em\allowbreak{} 可知

\[
\RB{134.129}{17.654}{15.210}{%
\RG{29.398}{13.639}{7.9701}{mi}{114}\RG{23.183}{10.367}{10.9092}{ex}{77}\RG{39.668}{8.839}{10.9092}{ex}{0}\RG{126.414}{8.839}{10.9092}{ex}{1}\RG{0.000}{0.003}{10.9091}{mi}{85}\RG{11.668}{0.003}{10.9091}{rm}{61}\RG{44.668}{0.003}{10.9091}{mi}{85}\RG{55.730}{0.003}{10.9091}{sy}{92}\RT{65.428}{0.003}{10.9091}{k}\RT{71.188}{0.003}{10.9091}{e}\RT{76.031}{0.003}{10.9091}{r}\RG{80.308}{0.003}{10.9091}{rm}{40}\RG{84.551}{0.003}{10.9091}{mi}{84}\RG{94.865}{0.003}{10.9091}{sy}{0}\RG{105.774}{0.003}{10.9091}{mi}{21}\RG{116.520}{0.003}{10.9091}{mi}{73}\RG{122.172}{0.003}{10.9091}{rm}{41}\RG{131.414}{0.003}{10.9091}{mi}{59}\RG{112.138}{-1.633}{7.9701}{mi}{106}\RG{24.074}{-13.084}{7.9701}{mi}{106}\RG{27.958}{-13.084}{7.9701}{rm}{61}\RG{34.542}{-13.084}{7.9701}{rm}{49}%
}
\]
从而\allowbreak{}\hskip.1em \RB{17.724}{8.684}{3.225}{%
\RG{0.000}{0.002}{10.9091}{mi}{84}\RG{7.890}{0.002}{10.9091}{sy}{106}\RG{10.920}{-1.687}{7.9701}{mi}{85}%
}\hskip.1em\allowbreak{} 有一组特征向量基，所述事实得证。

现设\allowbreak{}\hskip.1em \RB{11.889}{7.949}{0.502}{%
\RG{0.000}{-0.002}{10.9091}{mi}{77}%
}\hskip.1em\allowbreak{} 可对角化。由分块相似变换及上述事实，矩阵

\[
\RB{162.475}{8.687}{3.222}{%
\RG{0.000}{0.005}{10.9091}{mi}{67}\RG{11.608}{0.005}{10.9091}{rm}{61}\RG{23.123}{0.005}{10.9091}{mi}{65}\RG{33.729}{0.005}{10.9091}{rm}{43}\RG{44.638}{0.005}{10.9091}{mi}{102}\RG{51.153}{0.005}{10.9091}{rm}{40}\RG{55.395}{0.005}{10.9091}{mi}{65}\RG{63.578}{0.005}{10.9091}{rm}{41}\RG{67.821}{0.005}{10.9091}{mi}{59}\RG{94.472}{0.005}{10.9091}{mi}{68}\RG{106.852}{0.005}{10.9091}{rm}{61}\RG{118.368}{0.005}{10.9091}{mi}{65}\RG{128.974}{0.005}{10.9091}{sy}{0}\RG{139.883}{0.005}{10.9091}{mi}{102}\RG{146.398}{0.005}{10.9091}{rm}{40}\RG{150.640}{0.005}{10.9091}{mi}{65}\RG{158.822}{0.005}{10.9091}{rm}{41}%
}
\]
均可对角化，并且\allowbreak{}\hskip.1em \RB{50.566}{8.193}{0.738}{%
\RG{0.000}{0.002}{10.9091}{mi}{67}\RG{8.564}{0.002}{10.9091}{mi}{68}\RG{20.943}{0.002}{10.9091}{rm}{61}\RG{32.458}{0.002}{10.9091}{mi}{68}\RG{41.775}{0.002}{10.9091}{mi}{67}%
}\hskip.1em\allowbreak{} 。对\allowbreak{}\hskip.1em \RB{8.791}{8.193}{0.738}{%
\RG{0.000}{0.002}{10.9091}{mi}{67}%
}\hskip.1em\allowbreak{} 的任一特征子空间\allowbreak{}\hskip.1em \RB{23.198}{7.953}{2.372}{%
\RG{0.000}{0.002}{10.9091}{mi}{86}\RG{6.364}{-1.776}{7.9701}{mi}{21}\RG{14.832}{0.002}{10.9091}{rm}{61}%
}\hskip.1em\allowbreak{} ker\allowbreak{}\hskip.1em \RB{46.602}{8.684}{3.225}{%
\RG{0.000}{0.002}{10.9091}{rm}{40}\RG{4.243}{0.002}{10.9091}{mi}{67}\RG{15.244}{0.002}{10.9091}{sy}{0}\RG{26.153}{0.002}{10.9091}{mi}{21}\RG{32.513}{0.002}{10.9091}{mi}{73}\RG{37.312}{-1.635}{7.9701}{mi}{110}\RG{42.949}{0.002}{10.9091}{rm}{41}%
}\hskip.1em\allowbreak{} ，若\allowbreak{}\hskip.1em \RB{30.397}{7.953}{2.372}{%
\RG{0.000}{0.002}{10.9091}{mi}{122}\RG{8.583}{0.002}{10.9091}{sy}{50}\RG{18.886}{0.002}{10.9091}{mi}{86}\RG{25.250}{-1.776}{7.9701}{mi}{21}%
}\hskip.1em\allowbreak{} ，则

\[
\RB{116.936}{8.680}{3.229}{%
\RG{0.000}{-0.002}{10.9091}{mi}{67}\RG{8.578}{-0.002}{10.9091}{rm}{40}\RG{12.820}{-0.002}{10.9091}{mi}{68}\RG{22.137}{-0.002}{10.9091}{mi}{122}\RG{27.708}{-0.002}{10.9091}{rm}{41}\RG{34.963}{-0.002}{10.9091}{rm}{61}\RG{46.496}{-0.002}{10.9091}{mi}{68}\RG{55.831}{-0.002}{10.9091}{rm}{40}\RG{60.073}{-0.002}{10.9091}{mi}{67}\RG{68.637}{-0.002}{10.9091}{mi}{122}\RG{74.204}{-0.002}{10.9091}{rm}{41}\RG{81.458}{-0.002}{10.9091}{rm}{61}\RG{92.992}{-0.002}{10.9091}{mi}{21}\RG{99.352}{-0.002}{10.9091}{mi}{68}\RG{108.668}{-0.002}{10.9091}{mi}{122}\RG{114.221}{-0.002}{10.9091}{mi}{59}%
}
\]
所以\allowbreak{}\hskip.1em \RB{11.511}{7.956}{2.368}{%
\RG{0.000}{0.005}{10.9091}{mi}{86}\RG{6.364}{-1.772}{7.9701}{mi}{21}%
}\hskip.1em\allowbreak{} 是\allowbreak{}\hskip.1em \RB{9.271}{7.956}{0.500}{%
\RG{0.000}{0.005}{10.9091}{mi}{68}%
}\hskip.1em\allowbreak{} 的不变子空间。由上述事实，\allowbreak{}\hskip.1em \RB{21.778}{8.687}{3.661}{%
\RG{0.000}{0.005}{10.9091}{mi}{68}\RG{9.335}{0.005}{10.9091}{sy}{106}\RG{12.365}{-1.683}{7.9701}{mi}{86}\RG{17.297}{-3.089}{5.9776}{mi}{21}%
}\hskip.1em\allowbreak{} 可对角化。在每个\allowbreak{}\hskip.1em \RB{11.511}{7.956}{2.368}{%
\RG{0.000}{0.005}{10.9091}{mi}{86}\RG{6.364}{-1.772}{7.9701}{mi}{21}%
}\hskip.1em\allowbreak{} 中取\allowbreak{}\hskip.1em \RB{9.271}{7.956}{0.500}{%
\RG{0.000}{0.005}{10.9091}{mi}{68}%
}\hskip.1em\allowbreak{} 的特征向量基，并将这些基合并，便得到\allowbreak{}\hskip.1em \RB{22.067}{8.191}{2.606}{%
\RG{0.000}{-0.000}{10.9091}{mi}{67}\RG{7.953}{-0.000}{10.9091}{mi}{59}\RG{12.796}{-0.000}{10.9091}{mi}{68}%
}\hskip.1em\allowbreak{} 的一组公共特征向量基。因此

\[
\RB{69.714}{15.150}{7.979}{%
\RG{23.923}{7.384}{10.9091}{rm}{49}\RG{0.000}{0.004}{10.9091}{mi}{65}\RG{11.212}{0.004}{10.9091}{rm}{61}\RG{30.573}{0.004}{10.9091}{rm}{40}\RG{34.816}{0.004}{10.9091}{mi}{67}\RG{45.817}{0.004}{10.9091}{rm}{43}\RG{56.726}{0.004}{10.9091}{mi}{68}\RG{66.061}{0.004}{10.9091}{rm}{41}\RG{23.923}{-7.479}{10.9091}{rm}{50}\RR{23.923}{2.732}{5.455}{0.436}%
}
\]
在这组基下也是对角矩阵，即\allowbreak{}\hskip.1em \RB{8.366}{8.309}{0.502}{%
\RG{0.000}{-0.002}{10.9091}{mi}{65}%
}\hskip.1em\allowbreak{} 可对角化。\allowbreak{}\hskip.1em \RB{8.366}{7.993}{0.502}{%
\RG{0.000}{-0.002}{10.9091}{am}{3}%
}\hskip.1em\allowbreak{}

小结：先把对称分块矩阵分解为\allowbreak{}\hskip.1em \RB{44.353}{8.682}{3.227}{%
\RG{0.000}{0.000}{10.9091}{mi}{65}\RG{10.729}{0.000}{10.9091}{rm}{43}\RG{21.760}{0.000}{10.9091}{mi}{102}\RG{28.275}{0.000}{10.9091}{rm}{40}\RG{32.518}{0.000}{10.9091}{mi}{65}\RG{40.700}{0.000}{10.9091}{rm}{41}%
}\hskip.1em\allowbreak{} 与\allowbreak{}\hskip.1em \RB{44.352}{8.682}{3.227}{%
\RG{0.000}{0.000}{10.9091}{mi}{65}\RG{10.728}{0.000}{10.9091}{sy}{0}\RG{21.760}{0.000}{10.9091}{mi}{102}\RG{28.275}{0.000}{10.9091}{rm}{40}\RG{32.517}{0.000}{10.9091}{mi}{65}\RG{40.699}{0.000}{10.9091}{rm}{41}%
}\hskip.1em\allowbreak{} ，再利用它们可交换而构造公共特征向量基。

题目来源：谢启鸿《复旦大学高等代数习题集（2023版）》，第93页。


\end{document}
